\documentclass[11pt,a4paper]{article}

\usepackage[T1]{fontenc}
\usepackage[utf8]{inputenc}
\usepackage{lmodern}
\usepackage{amsmath,amssymb,amsthm,mathtools}
\usepackage{geometry}
\usepackage{microtype}
\usepackage{enumitem}
\usepackage{fancyhdr}
\usepackage{bm}

\geometry{margin=27mm}
\setlength{\parindent}{1.25em}
\setlength{\parskip}{0pt}
\setlength{\headheight}{14pt}
\allowdisplaybreaks

\pagestyle{fancy}
\fancyhf{}
\fancyhead[L]{Helffer--Sj\"ostrand, Multiple Wells I}
\fancyfoot[C]{\thepage}
\renewcommand{\headrulewidth}{0pt}

\newcommand{\R}{\mathbb R}
\newcommand{\C}{\mathbb C}
\newcommand{\N}{\mathbb N}
\newcommand{\Id}{\operatorname{Id}}
\newcommand{\Sp}{\operatorname{Sp}}
\newcommand{\dist}{\operatorname{dist}}
\newcommand{\supp}{\operatorname{supp}}
\newcommand{\Lip}{\operatorname{Lip}}
\newcommand{\Hol}{\operatorname{Hol}}
\newcommand{\Ran}{\operatorname{Ran}}
\newcommand{\Ker}{\operatorname{Ker}}
\newcommand{\Tr}{\operatorname{Tr}}
\newcommand{\cO}{\mathcal O}
\newcommand{\Ot}{\widetilde{\mathcal O}}
\newcommand{\eps}{\varepsilon}
\newcommand{\dd}{\,d}
\newcommand{\halfN}{\tfrac12\N}
\newcommand{\norm}[1]{\left\lVert #1\right\rVert}
\newcommand{\ip}[2]{\left(#1\mid #2\right)}
\newcommand{\dvec}{\overset{\rightarrow}{d}}
\newcommand{\cS}{\mathcal S}
\newcommand{\cL}{\mathcal L}
\newcommand{\cH}{\mathcal H}
\newcommand{\cE}{\mathcal E}
\newcommand{\cB}{\mathcal B}
\newcommand{\cF}{\mathcal F}
\newcommand{\vc}{\operatorname{v.c.}}
\newcommand{\Op}{\operatorname{Op}}

\numberwithin{equation}{section}
\newtheorem{theorem}{Theorem}[section]
\newtheorem{lemma}[theorem]{Lemma}
\newtheorem{proposition}[theorem]{Proposition}
\newtheorem{corollary}[theorem]{Corollary}
\theoremstyle{remark}
\newtheorem{remark}[theorem]{Remark}

\begin{document}

\begin{center}
{\LARGE\bfseries Multiple Wells in the Semi-Classical Limit I}\par
\vspace{1.2cm}
\begin{minipage}[t]{0.42\textwidth}\centering
{\large\itshape B. Helffer}\par\vspace{0.45cm}
D\'ept. de Math\'ematiques\\
Universit\'e de Nantes\\
2, chemin de la Houssini\`ere\\
F-44072 Nantes, FRANCE
\end{minipage}
\hfill
\begin{minipage}[t]{0.42\textwidth}\centering
{\large\itshape J. Sj\"ostrand}\par\vspace{0.45cm}
D\'ept. de Math\'ematiques\\
Universit\'e de Paris-Sud\\
F-91405 Orsay, FRANCE
\end{minipage}
\end{center}
\vspace{1.2cm}

\setcounter{section}{-1}
\section{Introduction}
This paper is devoted to the Schr\"odinger operator, but the general problem that we are interested in could be formulated also for a formally self-adjoint pseudodifferential operator
\begin{equation}
Pu(x,h)=(2\pi h)^{-n}\iint e^{\frac{i}{h}(x-y)\cdot\xi}P(x,\xi,h)u(y)\,dy\,d\xi,
\end{equation}
where
\begin{equation}
P(x,\xi,h)\sim\sum_{j=0}^{\infty}h^j p_j(x,\xi),\qquad h\to0.
\end{equation}
For large classes of such operators, Helffer and Robert [9], [10], [11], [12] have studied various types of asymptotics of eigenvalues. We are here interested in the behaviour of the eigenvalues tending to some given energy level $E_0\in\R$ when $h\to0$. Let $W$ be the energy ``surface'' given by $p_0(x,\xi)=E_0$. We assume that $W=\bigcup W_j$, where $W_j$ are disjoint and compact. These are the microlocal wells of the problem, to which the difficulties are localized, if one wants to study the eigenvalue asymptotics near $E_0$, either by the direct WKB--Maslov approach (see Maslov [17], Duistermaat [5], Candelpergher--Nosmas [3], Voros [26]) or by the trace formula approach combined with Fourier integral operators and functional calculus (see Helffer--Robert [9], [10], [11] and [12]).\footnote{Most of this work was done while the first author visited the ``Centre de Math\'ematiques, \'Ecole Polytechnique, Palaiseau''.}

In particular they proved under suitable assumptions that the union of the asymptotic eigenvalues that one can construct for each well gives the actual eigenvalues up to an error which is $\cO(h^\infty)$.

Our problem is to study much smaller corrections of size $\exp(-\mathrm{const.}/h)$ due to the interaction between such wells, for instance in situations where the sequences of asymptotic eigenvalues associated to the various wells are identical. From the microlocal point of view (see [24], for instance) it is natural to expect significant results only in the case of analytic pseudodifferential operators. The situation, however, is quite different in the second order case of the Schr\"odinger operator
\begin{equation}
P=-h^2\Delta+V(x),
\end{equation}
where $V$ is $C^\infty$, real valued and $\lim_{|x|\to\infty}V(x)>E_0$. The wells $U_1,\ldots,U_N$ are then different components of $\{x;\,V(x)\le E_0\}$ and $U_j$ is the $x$-space projection of $W_j$. The reason for this is that a normalized eigenfunction $u=u_h$ with $Pu=E_hu$, $E_h\to E_0$, $h\to0$, decays like $\exp(-\mathrm{const.}/h)$ on every compact set outside the wells. This has been proved in a slightly difficult context by Agmon [2] and in the semi-classical case by Simon [25]. From the point of view of [24] this is a special case of the method of non-characteristic deformations which in general should give microlocal exponential decay results outside the $W_j$ in the case of analytic pseudodifferential operators.

The classical example of double wells is the quartic oscillator
\[
P=-h^2\frac{d^2}{dx^2}+x^4-x^2
\]
in one dimension and has been discussed by Landau--Lifschitz [14], \S50, exercise no.~2, Reed--Simon [20] and Harrell [7]. In addition to this case our main source of inspiration has been the work of Harrell [8], which treats the case of symmetric double wells and gives a direct way of describing the splitting of eigenvalues. In trying to generalize Harrell's approach to the case of multiple wells, we encountered a mysterious square matrix of surface integrals of the same type as in [8].

Things became much clearer when we realized that this matrix is just the off-diagonal part (up to a very small error) of the matrix of $P$ in a suitable finite dimensional space and for a suitable basis, and finally our approach became quite different from [7], [8]. At this point we would like to mention also the interesting works of Davies [27], Mansour and M\"uller--Kristen [16], Jona-Lasinio, Martinelli, Scoppola [13], Simon [22], [23], Combes, Duclos, Seiler [29] and Damburg--Propin [28]. Simon [23] has obtained a general result on the splitting of the two lowest eigenvalues for double wells and we indicate at the end of this introduction the improvements given by our work.

In Section~1 we develop some weighted $L^2$ estimates in the spirit of Agmon [2], to be applied at various places in Sections~2 and~5. We also give some easy Hilbert space lemmas. In Section~2 we develop the main strategy, which is the following. Consider the Schr\"odinger operator (0.3) on $M$, where $M$ is $\R^n$ or a compact Riemannian manifold. In the first case we assume that $\lim_{|x|\to\infty}(V(x)-E_0)>0$. Let $U_1,\ldots,U_N$ be the wells for the given energy $E_0$. If $dx^2$ denotes the Riemannian metric on $M$, then the Agmon metric $\max(V(x)-E_0,0)dx^2$ gives a degenerate distance on $M$. We choose $M_j\subset M$ compact with smooth boundary such that $U_j\subset\mathring M_j$, $U_k\cap\partial M_j=\varnothing$ when $k\ne j$, and let $P_{M_j}$ be the self-adjoint realization of $P$ in $L^2(M_j)$ with Dirichlet boundary condition. Then for normalized eigenfunctions $u=u_h$, $P_{M_j}u=E_hu$, $E_h\to E_0$, $h\to0$, we have
\begin{equation}
\norm{u e^{d(\cdot,U_j)/h}}_{L^2(M_j)}+
\norm{\nabla\bigl(u e^{d(\cdot,U_j)/h}\bigr)}_{L^2(M_j)}
=\cO\!\left(e^{o(1)/h}\right),\qquad h\to0.
\end{equation}
Naturally we take $M_j$ as large as possible, so that
\[
B(U_j,S)=\{x\in M;\ d(x,U_j)\le S\}\subset M_j,
\]
where $S<S_0:=\min_{j\ne k}d(U_j,U_k)$ is close to $S_0$. We now consider a sequence of intervals $I(h)$, translated by $E_0$ if necessary, with $h\to0$, such that $P_{M_j}$ and $P$ have no eigenvalues outside $I(h)$ which are ``exponentially close'' to $I(h)$. Write
\[
\Sp(P_{M_j})\cap I(h)=\{\mu_{j,1},\ldots,\mu_{j,m_j}\},
\]
and let $\varphi_{j,k}$ be a corresponding orthonormal basis of eigenfunctions. We then prove, using only the energy decay and elementary arguments, that $I(h)$ contains exactly $m_1+\cdots+m_N$ eigenvalues of $P$. Moreover, if $\chi_j\in C_0^\infty(M_j)$ is equal to $1$ near $B(U_j,S)$ and $\psi_{j,k}=\chi_j\varphi_{j,k}$, if
\[
E=\langle\psi_{1,1},\ldots,\psi_{N,m_N}\rangle\subset L^2(M)
\]
and $F\subset L^2(M)$ is the space associated to $\Sp(P)\cap I(h)$, then $\dist(E,F)=\cO(e^{-S/h})$ in the sense defined in Section~1. We notice that $P|_F$ can be identified with $\tau P|_E$, where $\tau$ is the projection onto $E$ along $F^\perp$. Now $\tau$ can be approximated modulo $\Ot(e^{-S/h})$ by an explicit operator in terms of the $\psi_{j,k}$, and since $P\psi_{j,k}$ are at a distance $\Ot(e^{-S/h})$ from $F$, we get the matrix of $\tau P|_E$ in the basis $\psi_{j,k}$ up to an error $\Ot(e^{-2S/h})$ (Theorem~2.9). The expressions $w_{\alpha,\beta}$ in this matrix are further compared with certain surface integrals appearing in Harrell [8]. In a situation with many symmetries, for example, one may know that all the local eigenvalues $\mu_{j,k}$ agree, and asymptotic information about the $\varphi_{j,k}$ then gives asymptotic information about the relative positions of the true eigenvalues of $P$ in $I(h)$, although in this case they are all situated in an interval of length $\exp(-(S_0+o(1))/h)$. These arguments also show that $\Sp(P_{M_j})\cap I(h)$ is independent of the choice of $M_j$ modulo $\Ot(e^{-2S/h})$.

The value of such general results depends very much on how explicitly the $\varphi_{j,k}$, and also to some extent the $\psi_{j,k}$, can be computed, so further restrictive hypotheses on the potentials seem necessary. Many different classes are possible and in particular in the one-dimensional case much can be done. Rather than starting with the one-dimensional case we have preferred to look at the low eigenvalues in the case $E_0=0$, $V\ge0$, $U_j=\{x^j\}$, $V''(x^j)>0$. In Section~3 we construct the asymptotics of the first finitely many eigenvalues and the corresponding eigenfunctions associated to each well. In the recent work [22], B. Simon has obtained these asymptotics but his results on the asymptotics of the eigenfunctions are not precise enough for our purposes. In Section~4 we prove that in the real-analytic case these constructions give analytic symbols. In Section~5 we first show that the approximate eigenfunctions, which are of the form $a(x,h)\exp(-d(x,U_j)/h)$, approximate the actual eigenfunctions of $P_{M_j}$ closely, even when the exponential decrease factor is taken into account, in certain regions where $d(\cdot,U_j)$ is smooth. In the case of a $C^\infty$ potential, the required estimates go somewhat beyond Agmon's results [2].

We then assume that whenever $d(U_j,U_k)$ is minimal, i.e. equal to $S_0$, the wells $U_j$ and $U_k$ are linked by a unique minimal geodesic $\gamma_{j,k}\subset\Omega_j\cup\Omega_k$ (or possibly a finite number of such geodesics), and that $\gamma_{j,k}$ is a non-degenerate minimum in a very explicit way that we specify in the space of curves joining $x^j$ to $x^k$. Then, if all the $\mu_{j,k}$ have the same asymptotics (otherwise, just pass to a smaller interval $I(h)$) and $\mu_{j,k}=\mu_{j,k'}$ for all $j,k,k'$, the coefficients of the interaction matrix can be computed and they are of the form $b_{\alpha,\beta}(h)e^{-S_0/h}$ where
\[
b_{\alpha,\beta}(h)\sim h^{m(\alpha,\beta)}
\sum_{j\in\frac12\N} b_{\alpha,\beta}^{(j)}h^j.
\]
Moreover, they are analytic symbols in the real-analytic case and the half-degree terms disappear in the case of the principal eigenvalues.

In one or several forthcoming papers we shall study the implications of our results on the eigenvalue splitting in various situations with symmetries and for small perturbations of such situations. In [23], Simon has proved a result on the splitting of the two lowest eigenvalues in the case of two non-degenerate minima for the potential under the implicit condition that each of the two eigenfunctions is sufficiently present in both wells. He shows that the splitting is then $\exp(-(S_0+o(1))/h)$, $h\to0$. Our Proposition~2.16 gives this upper bound for the splitting in this case. Adding the geometric hypothesis mentioned above we get a much more precise result: the difference between the two lowest eigenvalues is
\[
\Delta E=\sqrt{4\widehat w_{1,2}^{\,2}+(\mu_1-\mu_2)^2}+\cO(e^{-2S/h}),
\]
where
\[
\widehat w_{1,2}=h^{1/2}b(h)e^{-S_0/h},
\qquad
b(h)\sim\sum_{j=0}^{\infty}b_jh^j,\quad b_0<0.
\]
This gives a precise lower bound for the splitting, and it follows from Section~6 (written partly in order to recover Simon's result completely) that this bound is still valid when the additional geometric condition is dropped.

We would like to thank D. Robert for several valuable discussions and for having indicated to us many interesting articles and preprints.

\section{Two preparations}
In this paper we will frequently need exponential decay estimates in the spirit of Agmon [2] (see also Carmona--Simon [4] for lower bounds). Here we first prove some a priori estimates that will make the proofs of the various decay results easy.

\begin{theorem}
Let $\Omega\subset\R^n$ be a bounded open domain with $C^2$ boundary, let $V\in C(\overline\Omega;\R)$ and $\phi\in\Lip(\overline\Omega;\R)$. Then $\nabla\phi$ is well defined in $L^\infty(\Omega)$ as the almost everywhere limit of $\nabla(\chi_\eps*\phi)=\chi_\eps*\nabla\phi$, $\eps\to0$, where $\chi_\eps\in C_0^\infty(B(0,\eps))$ is a standard mollifier. For every $u\in C^2(\overline\Omega;\R)$ with $u|_{\partial\Omega}=0$ we have
\begin{equation}
 h^2\left(\int_\Omega |\nabla(e^{\phi/h}u)|^2\,dx+
 \int_\Omega\bigl(V(x)-|\nabla\phi(x)|^2\bigr)e^{2\phi/h}u^2\,dx\right)
 =\int_\Omega e^{2\phi/h}Pu(x)u(x)\,dx,
\end{equation}
where $P=-h^2\Delta+V(x)$, $0<h\le1$.

In particular, if $V(x)-|\nabla\phi(x)|^2\ge0$ a.e. and $V(x)-|\nabla\phi(x)|^2=F(x)^2$, $0\le F\in L^\infty(\Omega)$, then
\begin{equation}
h^2\norm{\nabla(e^{\phi/h}u)}^2+\frac12\norm{Fe^{\phi/h}u}^2
\le \frac12\norm{F^{-1}e^{\phi/h}Pu}^2,
\end{equation}
where the norms are taken in $L^2(\Omega)$ (and the right hand side is defined as $+\infty$ when $F^{-1}e^{\phi/h}Pu\notin L^2(\Omega)$).

More generally, if $V(x)-|\nabla\phi(x)|^2=F_+(x)^2-F_-(x)^2$ a.e., where $0\le F_\pm\in L^\infty(\Omega)$, then
\begin{equation}
h^2\norm{\nabla(e^{\phi/h}u)}^2+\frac12\norm{F_+e^{\phi/h}u}^2
\le \norm{\frac{e^{\phi/h}Pu}{F_++F_-}}^2
+\frac32\norm{F_-e^{\phi/h}u}^2.
\end{equation}
\end{theorem}

\begin{proof}
Let first $\phi\in C^2(\overline\Omega)$. Then by Green's formula
\begin{align*}
h^2\int_\Omega |\nabla(e^{\phi/h}u)|^2\,dx
&=-\int_\Omega h^2\Delta(e^{\phi/h}u)e^{\phi/h}u\,dx\\
&=-\int_\Omega e^{2\phi/h}h^2\Delta u\,u\,dx
-2\int_\Omega e^{2\phi/h}\sum_j\frac{\partial\phi}{\partial x_j}h\frac{\partial u}{\partial x_j}u\,dx\\
&\qquad-\int_\Omega h^2\Delta(e^{\phi/h})u e^{\phi/h}u\,dx
=:I+II+III.
\end{align*}
Since $u^2$ vanishes to the second order on the boundary,
\begin{align*}
II
&=-\sum_j\int_\Omega e^{2\phi/h}\frac{\partial\phi}{\partial x_j}h\frac{\partial}{\partial x_j}(u^2)\,dx\\
&=\sum_j\int_\Omega h\frac{\partial}{\partial x_j}\left(e^{2\phi/h}\frac{\partial\phi}{\partial x_j}\right)u^2\,dx\\
&=\int_\Omega e^{2\phi/h}\left(2\sum_j\left(\frac{\partial\phi}{\partial x_j}\right)^2+h\sum_j\frac{\partial^2\phi}{\partial x_j^2}\right)u^2\,dx.
\end{align*}
On the other hand in $III$ we substitute
\[
-h^2\Delta(e^{\phi/h})=-\sum_j e^{\phi/h}\left[\left(\frac{\partial\phi}{\partial x_j}\right)^2+h\frac{\partial^2\phi}{\partial x_j^2}\right],
\]
so
\[
h^2\int_\Omega|\nabla(e^{\phi/h}u)|^2\,dx
=-\int_\Omega e^{2\phi/h}h^2\Delta u\,u\,dx+
\int_\Omega e^{2\phi/h}|\nabla\phi|^2u^2\,dx,
\]
and (1.1) follows.

Now let $\phi\in\Lip(\Omega)$ and extend $\phi$ to a Lipschitz function on $\R^n$ with compact support. Put $\phi_\eps=\chi_\eps*\phi$. Then $\nabla\phi_\eps$ is a bounded family in $L^\infty$, and if $\nabla\phi$ is the gradient in the distribution sense, then for every $\varphi\in C_0^\infty$,
\[
\left|\int\nabla\phi\,\varphi\,dx\right|
=\left|\lim_{\eps\to0}\int\nabla\phi_\eps\,\varphi\,dx\right|
\le C\norm{\varphi}_{L^1}.
\]
Thus $\nabla\phi\in L^\infty$. Moreover $\nabla\phi_\eps=\chi_\eps*\nabla\phi\to\nabla\phi$ a.e. as $\eps\to0$. Since (1.1) holds with $\phi$ replaced by $\phi_\eps$, the general case follows by dominated convergence.

To prove (1.2), estimate the right hand side of (1.1) by
\[
\norm{F^{-1}e^{\phi/h}Pu}\,\norm{Fe^{\phi/h}u}
\le\frac12\norm{F^{-1}e^{\phi/h}Pu}^2+\frac12\norm{Fe^{\phi/h}u}^2.
\]
To prove (1.3), with $F=F_++F_-$, use
\[
\norm{F^{-1}e^{\phi/h}Pu}^2+\frac14\norm{Fe^{\phi/h}u}^2
\le \norm{F^{-1}e^{\phi/h}Pu}^2+\frac12\norm{F_+e^{\phi/h}u}^2+\frac12\norm{F_-e^{\phi/h}u}^2,
\]
and notice that
\[
\int_\Omega(V-|\nabla\phi|^2)e^{2\phi/h}u^2\,dx
=\norm{F_+e^{\phi/h}u}^2-\norm{F_-e^{\phi/h}u}^2.
\]
\end{proof}

\begin{remark}
Let $M$ be a compact Riemannian manifold. Then $dx$ will denote the induced volume element and $\nabla$ the exterior differential on functions. Since we have an induced metric on $T_x^*M$ for every $x\in M$, the adjoint $\nabla^*$ is a well defined differential operator from $1$-forms to functions. Locally, if we choose an orthonormal basis of $C^\infty$ vector fields $X_1,\ldots,X_n$, then
\[
\nabla u=\sum_{j=1}^n X_ju\,e_j,
\]
where $e_1,\ldots,e_n$ is the dual basis of $1$-forms. If $\omega=\sum\omega_je_j$ is a $1$-form, then
\[
\nabla^*\omega=\sum_jX_j^*(\omega_j),
\]
where $X_j^*=-X_j+g_j$, with $g_j$ a real $C^\infty$ function. If $f$ is a function, then $\nabla^*(f\omega)=f\nabla^*\omega-\nabla f\cdot\omega$. The Laplacian on functions is defined as $\Delta=-\nabla^*\nabla$. Theorem~1.1 still holds with $\R^n$ replaced by $M$, with the obvious modifications.
\end{remark}

As our second preparation, we discuss the distance between closed subspaces of a Hilbert space. Let $H$ be a Hilbert space and let $E,F$ be closed subspaces. Let $\Pi_E,\Pi_F$ be the orthogonal projections on $E$ and $F$ respectively. We define the non-symmetric distance by
\[
\dvec(E,F)=\sup_{x\in E,\ \norm{x}=1}d(x,F)
=\norm{(I-\Pi_F)|_E}=\norm{\Pi_E-\Pi_F\Pi_E}.
\]
Note that $\dvec(E,F)=0$ if and only if $E\subset F$, and that $\dvec(E,G)\le\dvec(E,F)+\dvec(F,G)$ if $G$ is a third closed subspace of $H$.

\begin{lemma}
If $\dvec(E,F)<1$, then $\Pi_F|_E:E\to F$ is injective and $\Pi_E|_F:F\to E$ has a bounded right inverse.
\end{lemma}
\begin{proof}
If $0\ne x\in E$ and $\Pi_Fx=0$, then $\dvec(E,F)=1$, so the injectivity is clear. On the other hand, if $x\in E$, we look for $y=\Pi_Fz$, $z\in E$, such that $x=\Pi_Ey=\Pi_E\Pi_Fz$. Writing this as $x=(I-(I-\Pi_E\Pi_F))z$, there is a unique solution with $\norm z\le C\norm x$, since $\norm{(I-\Pi_E\Pi_F)|_E}<1$.
\end{proof}

\begin{proposition}
If $\dvec(E,F)<1$ and $\dvec(F,E)<1$, then $\Pi_E|_F:F\to E$ and $\Pi_F|_E:E\to F$ are bijective with bounded inverses. Moreover $\dvec(E,F)=\dvec(F,E)$.
\end{proposition}
\begin{proof}
Only the last statement remains to prove. Clearly we may represent
\[
F=\{x+Ax;\ x\in E\},\qquad A\in\mathcal L(E,E^\perp).
\]
Then $\Pi_E(x+Ax)=x$, and hence
\[
\dvec(F,E)^2=\sup_{x\in E\setminus\{0\}}\frac{\norm{Ax}^2}{\norm{x+Ax}^2}
=\sup_{x\in E\setminus\{0\}}\frac{\norm{Ax}^2}{\norm{x}^2+\norm{Ax}^2}.
\]
For $x\in E$, write $\Pi_Fx=y+Ay$, $y\in E$. The condition
\[
\ip{y+Ay}{t+At}=\ip{x}{t+At},\qquad t\in E,
\]
gives $y+A^*Ay=x$, hence $y=(I+A^*A)^{-1}x$ and
\[
\Pi_Fx=(I+A^*A)^{-1}x+A(I+A^*A)^{-1}x.
\]
Since $(I+A^*A)^{-1}-I=-A^*A(I+A^*A)^{-1}$,
\[
\Pi_Fx-x=-A^*A(I+A^*A)^{-1}x+A(I+A^*A)^{-1}x,
\]
so
\begin{align*}
\norm{\Pi_Fx-x}^2
&=\norm{A^*A(I+A^*A)^{-1}x}^2
 +\norm{A(I+A^*A)^{-1}x}^2.
\end{align*}
Putting $x=(I+A^*A)^{1/2}y$ gives
\begin{align*}
\dvec(E,F)^2
&=\sup_{y\ne0}(\norm y^2+\norm{Ay}^2)^{-1}\norm{Ay}^2\\
&=\dvec(F,E)^2.
\end{align*}
\end{proof}

\section{General results about the interaction between wells}
In the following $M$ will denote $\R^n$ or a connected $n$-dimensional compact Riemannian manifold. Let $V\in C^\infty(M;\R)$ and put
\[
P=-h^2\Delta+V.
\]
This is a self-adjoint operator on $L^2(M)$ when $M$ is compact and has a discrete spectrum of finite multiplicity. Our goal is to study the behaviour of the eigenvalues of $P$ near some fixed energy level $E=E_0$ when $h\to0$. Changing $V$ we may suppose that $E_0=0$. In the case $M=\R^n$ we make the following assumption, which is automatic when $M$ is compact:
\begin{equation}
\text{There is an interval }I_0=(-\infty,a_0],\ a_0>0,\text{ such that }V^{-1}(I_0)\text{ is compact.}
\end{equation}
Then $P$ is a semibounded symmetric operator on $C_0^\infty(M)$ and we denote by $P$ also its Friedrichs extension. The following property is clear when $M$ is compact and follows in the case $M=\R^n$ from the max--min principle:
\begin{equation}
\begin{minipage}{0.88\textwidth}
For each $I_1=(-\infty,a_1]$, $a_1\le a_0$, the spectrum $\Sp(P)\cap I_1$ is discrete, finite, and of finite multiplicity. Moreover there is a constant $N_0\ge0$ such that the number of eigenvalues of $P$ in $I_1$, counted with multiplicity, is at most $h^{-N_0}$ for $0<h\le1/N_0$.
\end{minipage}
\end{equation}

Having fixed the energy level to $E_0=0$, decompose
\[
\{x\in M;\ V(x)\le0\}=U_1\cup\cdots\cup U_N
\]
into finitely many compact disjoint sets. The $U_j$ are called wells. In $M$ we consider the Agmon metric $\max(V,0)dx^2$, where $dx^2$ denotes the Riemannian metric on $M$. Let $d(x,y)$ be the associated degenerate distance and put
\[
S_0=\min_{j\ne k}d(U_j,U_k).
\]
Fix $S\in(0,S_0)$ and put $B(U_j,S)=\{x\in M;\ d(x,U_j)<S\}$. We can find compact manifolds with $C^2$ boundary $M_j\Subset M$ such that
\begin{equation}
B(U_j,S)\subset\mathring M_j,\qquad M_j\cap U_k=\varnothing\quad(k\ne j).
\end{equation}
On $L^2(M_j)$ we consider the self-adjoint Dirichlet realization $P_{M_j}$ of $P$. From now on $h$ need not vary in a full interval $(0,\eps_0]$, but only in a subset $J$ having $0$ as an accumulation point.

Let $I=[\alpha(h),\beta(h)]$, $h\in J$, with $\alpha(h),\beta(h)\to0$ as $h\to0$, and assume that there is a function $a(h)>0$ such that
\begin{equation}
|\log a(h)|=o(1/h),\qquad h\to0,
\end{equation}
and
\begin{equation}
P,P_{M_1},\ldots,P_{M_N}\text{ have no spectrum in }(\alpha-2a,\alpha)\text{ or }(\beta,\beta+2a).
\end{equation}
Write
\[
\Sp(P)\cap I=\{\lambda_1,\ldots,\lambda_M\},\qquad
\Sp(P_{M_j})\cap I=\{\mu_{j,1},\ldots,\mu_{j,m_j}\},
\]
and let $u_1,\ldots,u_M\in L^2(M)$ and $\varphi_{j,1},\ldots,\varphi_{j,m_j}\in L^2(M_j)$ be corresponding orthonormal systems. There is a constant $N_0$ such that
\[
M+m_1+\cdots+m_N=\cO(h^{-N_0}).
\]

\begin{remark}
All our estimates are uniform with respect to additional parameters if the assumptions satisfy suitable uniformity conditions.
\end{remark}

\begin{remark}
We shall see below that the eigenvalues $\mu_{j,k}$ are modified only by a term $\cO(e^{-2S/h})$ when the choice of $M_j$ is modified. Hence the assumptions above do not really depend on the domains $M_j$. It is plausible that $M_j=M\setminus\bigcup_{k\ne j}U_k$ would be an optimal choice, but the advantages seem small compared with the additional technical difficulties.
\end{remark}

\begin{lemma}
Let $\phi_j(x)=d(x,U_j)$, $x\in M_j$. Then for every $\delta>0$,
\begin{equation}
\norm{\nabla(e^{\phi_j/h}\varphi_{j,k})}+\norm{e^{\phi_j/h}\varphi_{j,k}}
=\cO(e^{\delta/h}),\qquad h\to0.
\end{equation}
\end{lemma}
\begin{proof}
We have $V(x)-|\nabla\phi_j|^2\ge0$ in $M_j\setminus U_j$. If $\phi=(1-\delta)\phi_j$, then
\[
V(x)-|\nabla\phi|^2\ge(2\delta-\delta^2)V(x)
\]
in $M_j\setminus U_j$. We apply (1.3) with $u=\varphi_{j,k}$ and with the operator $-h^2\Delta+(V-\mu_{j,k})$. For $h$ sufficiently small we can choose $F_+$ strictly positive and equal to $\sqrt{V-\mu_{j,k}-|\nabla\phi|^2}$ when $d(x,U_j)\ge\delta$, while $F_-$ is supported in $B(U_j,\delta)$. Then
\[
h^2\norm{\nabla(e^{\phi/h}\varphi_{j,k})}^2+
\norm{e^{\phi/h}\varphi_{j,k}}^2\le C_\delta e^{2\delta/h}.
\]
This also controls $\norm{e^{\phi/h}\nabla\varphi_{j,k}}$ and gives the result.
\end{proof}

Let $\chi_j\in C_0^\infty(\mathring M_j)$ be equal to $1$ in a neighborhood of $B(U_j,S)$ and put $\psi_{j,k}=\chi_j\varphi_{j,k}$. Then there is $S_1\in(S,S_0)$ such that
\begin{equation}
P\psi_{j,k}=\mu_{j,k}\psi_{j,k}+\cO(e^{-S_1/h})\quad\text{in }L^2(M_j).
\end{equation}
Let $E_j\subset L^2(M)$ be the subspace spanned by $\psi_{j,1},\ldots,\psi_{j,m_j}$, let $E=\bigoplus_{j=1}^N E_j$, and let $F$ be the space spanned by $u_1,\ldots,u_M$.

\begin{theorem}
For every $\sigma<S_1$,
\[
\dvec(E,F),\ \dvec(F,E)=\cO(e^{-\sigma/h}).
\]
Moreover, for $h$ sufficiently small there is a bijection
\[
b:\Sp(P)\cap I\longrightarrow\bigcup_{j=1}^N(\Sp(P_{M_j})\cap I)
\]
such that $b(\lambda)-\lambda=\cO(e^{-\sigma/h})$ for every $\sigma<S_1$.
\end{theorem}
The main goal of this paper is to improve the second part of the theorem. We shall prove it by a perturbation argument.

\begin{proposition}
Let $A$ be a self-adjoint operator in a Hilbert space $H$. Let $I\subset\R$ be a compact interval and let $\psi_1,\ldots,\psi_N\in H$ be linearly independent, with $\mu_1,\ldots,\mu_N\in I$, such that
\begin{equation}
A\psi_j=\mu_j\psi_j+r_j,\qquad\norm{r_j}\le\eps.
\end{equation}
Let $a>0$ and assume that
\[
\Sp(A)\cap\bigl((I+B(0,2a))\setminus I\bigr)=\varnothing.
\]
If $E=\langle\psi_1,\ldots,\psi_N\rangle$ and $F$ is the spectral subspace associated with $\Sp(A)\cap I$, then
\begin{equation}
\dvec(E,F)\le \frac{N^{1/2}\eps}{a\,\lambda_{\min}(S)^{1/2}},
\end{equation}
where $S=(\ip{\psi_j}{\psi_k})$.
\end{proposition}

\begin{remark}
If $\Sp(A)\cap I$ is discrete of finite multiplicity and the right hand side of (2.9) is strictly smaller than $1$, then $A$ has at least $N$ eigenvalues in $I$.
\end{remark}

\begin{proof}[Proof of Proposition 2.5]
Let $\lambda\in\C\setminus(\{\mu_1,\ldots,\mu_N\}\cup\Sp(A))$. By (2.8),
\begin{equation}
(A-\lambda)^{-1}\psi_j=(\mu_j-\lambda)^{-1}\psi_j
-(\mu_j-\lambda)^{-1}(A-\lambda)^{-1}r_j.
\end{equation}
If $\gamma_R$ is the oriented boundary of $(I+B(0,a))\times i[-R,R]$, then
\[
\Pi_F\psi_j=\frac1{2\pi i}\int_{\gamma_R}(\mu_j-\lambda)^{-1}\psi_j\,d\lambda
-\frac1{2\pi i}\int_{\gamma_R}(\mu_j-\lambda)^{-1}(A-\lambda)^{-1}r_j\,d\lambda.
\]
The first integral equals $\psi_j$ and the second tends, as $R\to\infty$, to the sum of the two vertical contour integrals at $\beta+a$ and $\alpha-a$, where $I=[\alpha,\beta]$. On either vertical line,
\[
\norm{(\mu_j-\lambda)^{-1}(A-\lambda)^{-1}r_j}\le\frac{\eps}{a^2+t^2}.
\]
Hence
\[
\norm{\Pi_F\psi_j-\psi_j}\le\frac\eps\pi\int_{-\infty}^{\infty}\frac{dt}{a^2+t^2}=\frac\eps a.
\]
If $u=\sum_{j=1}^N\alpha_j\psi_j\in E$, then
\[
\norm u^2=(S\alpha\mid\alpha)\ge\lambda_{\min}(S)\norm\alpha^2,
\]
and therefore
\[
\norm{\Pi_Fu-u}\le\sum_j|\alpha_j|\norm{\Pi_F\psi_j-\psi_j}
\le\frac{N^{1/2}\eps}{a\lambda_{\min}(S)^{1/2}}\norm u.
\]
\end{proof}

\begin{proof}[Proof of Theorem 2.4]
To estimate $\dvec(E,F)$ we apply Proposition~2.5 with $A=P$ and the vectors $\psi_{j,k}$. Equation (2.7) gives $\eps=\cO(e^{-S_1/h})$, the number of approximate eigenvalues is $\cO(h^{-N_0})$, and
\[
\ip{\psi_{j,k}}{\psi_{\ell,m}}=\delta_{j\ell}\delta_{km}+\cO(e^{-\sigma/h})
\]
for every $\sigma<S_0$. Thus $\lambda_{\min}(S)=1+\cO(e^{-\sigma/h})$ and the first estimate follows.

To estimate $\dvec(F,E)$, let $u=u_k$ and $\lambda=\lambda_k$. We use the following decay lemma.
\end{proof}

\begin{lemma}
Let $\widetilde\phi(x)=\min_j d(x,U_j)$. Then for every $\eps>0$,
\begin{equation}
\norm{\nabla(e^{(1-\eps)\widetilde\phi/h}u)}+
\norm{e^{(1-\eps)\widetilde\phi/h}u}=\cO(e^{\eps/h}),\qquad h\to0.
\end{equation}
\end{lemma}
\begin{proof}
When $M$ is compact the proof is a minor modification of Lemma~2.3. When $M=\R^n$, the Friedrichs realization gives
\[
h^2\norm{\nabla u}^2+\int V(x)|u(x)|^2\,dx=\lambda\norm u^2.
\]
For small $h$, $|\lambda|\le1$, and if $\min V=-b$ then $h^2\norm{\nabla u}^2\le b+1$. By approximating $u$ with $\chi(x/R)u(x)$ one obtains (1.1) and (1.3) on $\R^n$ for weights constant at infinity. Let $\chi_R(t)=t$ on $[0,R]$, $\chi_R(t)=R$ for $t>R$, and put
\[
\phi_R(x)=(1-\eps)\chi_R(\widetilde\phi(x)).
\]
Outside the wells,
\[
V(x)-\lambda-|\nabla\phi_R|^2\ge(2\eps-\eps^2)V(x)-\lambda.
\]
Apply (1.3), taking $F_+$ and $F_-$ as the square roots of the positive and negative parts of this expression. One obtains
\[
\begin{split}
h^2\norm{\nabla(e^{\phi_R/h}u)}^2
&+\int_{(2\eps-\eps^2)V-\lambda\ge0}
((2\eps-\eps^2)V-\lambda)e^{2\phi_R/h}|u|^2\,dx\\
&\le C\int_{(2\eps-\eps^2)V-\lambda\le0}
e^{2(1-\eps)\widetilde\phi/h}|u|^2\,dx.
\end{split}
\]
The right hand side is independent of $R$ and is $\cO(e^{o(1)/h})$ as $\eps\to0$. Letting $R\to\infty$ gives (2.11). In the compact case the factor $1-\eps$ can be replaced by $1$.
\end{proof}

Choose $\widehat\chi_j\in C_0^\infty(\mathring M_j)$, $j=1,\ldots,N$, equal to $1$ near $U_j$ and with disjoint supports, and define $\widehat\chi_0\in C_0^\infty(M)$ so that $1=\sum_{j=0}^N\widehat\chi_j$. Lemma~2.7 gives a constant $C_0>0$ such that
\begin{align}
\widehat\chi_0u&=\cO(e^{-C_0/h})\quad\text{in }L^2(M),\tag{2.12}\\
P\widehat\chi_ju&=\lambda\widehat\chi_ju+\cO(e^{-C_0/h})\quad\text{in }L^2(M).\tag{2.13}
\end{align}
Applying Proposition~2.5 with $A=P_{M_j}$ gives, for every $C_1<C_0$,
\begin{align}
\widehat\chi_ju&=\sum_{k=1}^{m_j}a_{j,k}\varphi_{j,k}+\cO(e^{-C_1/h}),\tag{2.14}\\
\sum_k|a_{j,k}|^2&\le\norm{\widehat\chi_ju}^2\le1.\tag{2.15}
\end{align}
After possibly decreasing $C_1$,
\begin{equation}
\widehat\chi_ju=\sum_{k=1}^{m_j}a_{j,k}\psi_{j,k}+\cO(e^{-C_1/h}),
\tag{2.16}
\end{equation}
and hence
\begin{equation}
u=\sum_{j=1}^N\sum_{k=1}^{m_j}a_{j,k}\psi_{j,k}+\cO(e^{-C_2/h}),\qquad0<C_2<C_1.
\tag{2.17}
\end{equation}
Since the vectors $u$ form an orthonormal basis for $F$, $\dvec(F,E)=\cO(e^{-C_3/h})$ for $0<C_3<C_2$. Proposition~1.4 then gives
\[
\dvec(E,F)=\dvec(F,E)=\cO(e^{-\sigma/h}),\qquad \sigma<S_1.
\]
Moreover, for small $h$, $M=m_1+\cdots+m_N$. To prove the spectral bijection, choose $\sigma<\sigma'<S_1$ and put $a=e^{-\sigma'/h}$. Cover the combined $\lambda$- and $\mu$-sets by disjoint intervals $I_1,\ldots,I_\ell$ separated by at least $2a$ and of length $\cO(h^{-N_0}a)$. Remark~2.6 applied to each interval shows that the numbers of $\lambda$'s and $\mu$'s agree in every $I_j$. Choosing a bijection within each interval gives
\[
b(\lambda)-\lambda=\cO(h^{-N_0}e^{-\sigma'/h})=\cO(e^{-\sigma/h}),
\]
which completes the proof of Theorem~2.4.

Keeping the hypotheses of Theorem~2.4, write $\alpha=(j,k)$ and $j(\alpha)=j$. Let $\Pi_0$ be the projection onto $E$ along $F^\perp$ in $L^2(M)$.

\begin{lemma}
For every $\sigma<S_1$,
\[
\norm{\Pi_0-\Pi_E}=\cO(e^{-\sigma/h}).
\]
\end{lemma}
\begin{proof}
For small $h$ we have $F=\{x+Ax;\ x\in E\}$, where $A:E\to E^\perp$ and $\norm A=\cO(e^{-\sigma/h})$. Then $F^\perp=\{y-A^*y;\ y\in E^\perp\}$. If $z=x+y$ with $x\in E$, $y\in E^\perp$, then $\Pi_Ez=x$ while $\Pi_0z=\widehat x$ is determined by $\widehat x-(x+y)=A^*\widehat y-\widehat y$. Thus $\widehat y=y$ and $\widehat x-x=A^*y$, proving the estimate.
\end{proof}

We have $\Pi_0=\Pi_0\Pi_F$, and the inverse of $\Pi_0|_F:F\to E$ is $\Pi_F|_E:E\to F$. Since $F$ and $F^\perp$ are stable under $P$,
\begin{equation}
\Pi_0P\Pi_Fu=\Pi_0Pu,\qquad u\in E.
\end{equation}
Thus, under this identification, $P|_F$ becomes $\Pi_0P|_E$.

\begin{theorem}
For every $\sigma<S_1$, the matrix of $\Pi_0P|_E$ in the basis $(\psi_\alpha)_\alpha$ is
\[
\operatorname{diag}(\mu_{1,1},\ldots,\mu_{N,m_N})+(\omega_{\alpha\beta})+
\cO(e^{-2\sigma/h}),
\]
where
\[
\omega_{\alpha\beta}=h^2\int
\chi_{j(\alpha)}\bigl(\varphi_\beta\nabla\varphi_\alpha-
\varphi_\alpha\nabla\varphi_\beta\bigr)\cdot\nabla\chi_{j(\beta)}\,dx.
\]
Moreover $\omega_{\alpha\beta}=\cO(e^{-\sigma/h})$ for every $\sigma<S_1$, and $\omega_{\alpha\beta}=\cO(e^{-2\sigma/h})$ when $j(\alpha)=j(\beta)$.
\end{theorem}
\begin{proof}
Write $S=(\ip{\psi_\alpha}{\psi_\beta})$. The orthogonal projection onto $E$ has the form
\[
\Pi_Eu=\sum_{\alpha,\beta}c_{\alpha\beta}\ip{u}{\psi_\beta}\psi_\alpha,
\qquad C=(c_{\alpha\beta})={}^tS^{-1}.
\]
Since $S=I+\cO(e^{-\sigma/h})$, the operator
\[
\tau u=\sum_\alpha\ip{u}{\psi_\alpha}\psi_\alpha
\]
satisfies
\begin{equation}
\norm{\tau-\Pi_0}=\cO(e^{-\sigma/h}),\qquad\sigma<S_1.
\end{equation}
Also
\[
P\psi_\beta=\mu_\beta\psi_\beta-h^2\left(2\nabla\chi_{j(\beta)}\cdot\nabla\varphi_\beta+
(\Delta\chi_{j(\beta)})\varphi_\beta\right).
\]
Hence
\begin{equation}
\Pi_0P\psi_\beta=\mu_\beta\psi_\beta-h^2\tau\left(2\nabla\chi_{j(\beta)}\cdot\nabla\varphi_\beta+
(\Delta\chi_{j(\beta)})\varphi_\beta\right)+\cO(e^{-2\sigma/h}).
\end{equation}
After inserting $\tau$ and integrating by parts, the terms with $\nabla\chi_{j(\alpha)}\cdot\nabla\chi_{j(\beta)}$ cancel and one obtains
\begin{equation}
\Pi_0P\psi_\beta=\mu_\beta\psi_\beta+
 h^2\sum_\alpha\left\{\int\chi_{j(\alpha)}
(\varphi_\beta\nabla\varphi_\alpha-
\varphi_\alpha\nabla\varphi_\beta)\cdot\nabla\chi_{j(\beta)}\,dx\right\}\psi_\alpha
+\cO(e^{-2\sigma/h}).
\end{equation}
\end{proof}

We notice that $\Pi_0P|_E=\Pi_EP|_E+\cO(e^{-2\sigma/h})$ and $\Pi_EP|_E$ is self-adjoint. The matrix $(w_{\alpha\beta})$ need not be symmetric because the $\psi_\alpha$ are not orthonormal. Integration by parts gives
\begin{equation}
w_{\alpha\beta}=(\mu_\alpha-\mu_\beta)t_{\alpha\beta}+w_{\beta\alpha},
\tag{2.22}
\end{equation}
where $S=I+T$, $T=(t_{\alpha\beta})$. If $\boldsymbol\psi=(\psi_{1,1},\ldots,\psi_{N,m_N})$, then $\mathbf e=\boldsymbol\psi S^{-1/2}$ is an orthonormal basis of $E$. In this basis
\[
(I+T)^{1/2}\bigl((\mu_\alpha\delta_{\alpha\beta})+(w_{\alpha\beta})\bigr)(I+T)^{-1/2}
\equiv(\mu_\alpha\delta_{\alpha\beta}+\widehat w_{\alpha\beta}),
\]
modulo $\cO(e^{-2\sigma/h})$, where
\[
\widehat w_{\alpha\beta}=w_{\alpha\beta}+\frac12t_{\alpha\beta}(\mu_\beta-\mu_\alpha)
=\frac12(w_{\alpha\beta}+w_{\beta\alpha}).
\]

\begin{theorem}
Let $S=(\ip{\psi_\alpha}{\psi_\beta})$ and let $\mathbf e=\boldsymbol\psi S^{-1/2}$ be the orthonormalization of the basis $\boldsymbol\psi$. For every $\sigma<S_1$, the matrix of $\Pi_0P|_E$ in this basis is
\begin{equation}
\operatorname{diag}(\mu_{1,1},\ldots,\mu_{N,m_N})+(\widehat w_{\alpha\beta})+
\cO(e^{-2\sigma/h}).
\end{equation}
Here
\[
\widehat w_{\alpha\beta}=\frac12(w_{\alpha\beta}+w_{\beta\alpha})
\]
and therefore
\[
\widehat w_{\alpha\beta}=\frac{h^2}{2}\int\Bigl[
\chi_{j(\alpha)}(\varphi_\beta\nabla\varphi_\alpha-
\varphi_\alpha\nabla\varphi_\beta)\cdot\nabla\chi_{j(\beta)}
+\chi_{j(\beta)}(\varphi_\alpha\nabla\varphi_\beta-
\varphi_\beta\nabla\varphi_\alpha)\cdot\nabla\chi_{j(\alpha)}\Bigr]dx.
\]
\end{theorem}

\begin{remark}
Let $\mathcal A=\{(j,k);\ 1\le j\le N,\ 1\le k\le m_j\}$ and write $\mathbf e=(e_\alpha)_{\alpha\in\mathcal A}$. The preceding theorem also gives the matrix of $P|_F$ in the basis $f_\alpha=\Pi_F(e_\alpha)$. Since $e_\alpha=f_\alpha+r_\alpha$, $r_\alpha\in F^\perp$, $\norm{r_\alpha}=\cO(e^{-\sigma/h})$, we have
\[
\ip{f_\alpha}{f_\beta}=\delta_{\alpha\beta}+\cO(e^{-2\sigma/h}).
\]
After orthonormalizing $(f_\alpha)$, the matrix of $P|_F$ is still of the form (2.23).
\end{remark}

\begin{theorem}
For $h$ sufficiently small there is a bijection
\[
b:\Sp(P|_F)\longrightarrow\Sp(\mu_\alpha\delta_{\alpha\beta}+\widehat w_{\alpha\beta})
\]
such that
\[
|b(\lambda)-\lambda|=\cO(e^{-2\sigma/h}),\qquad\sigma<S_1.
\]
\end{theorem}
For the lowest eigenvalues, the formal L.C.A.O. method using the $\psi_\alpha$ as approximate eigenfunctions gives precisely the matrix (2.23). Thus the result gives a rigorous form of the L.C.A.O. approximation.

\begin{remark}
Naturally the results are most interesting when $\mu_\alpha-\mu_\beta$ is small. Assume $0<a\le S_0$ and $S_0+a<2S$, and suppose that for every $\eps>0$,
\begin{equation}
\mu_\alpha-\mu_\beta=\cO\!\left(e^{-(a-\eps)/h}\right).
\end{equation}
Then, by (2.22),
\[
w_{\alpha\beta}=w_{\beta\alpha}\quad\bmod\ \cO\!\left(e^{-(S_0+a-\eps)/h}\right).
\]
Modulo errors of the same size, $w_{\alpha\beta}$ does not depend on the choices of $\chi_{j(\alpha)}$ and $\chi_{j(\beta)}$. If $d(\mathring U_{j(\alpha)},\mathring U_{j(\beta)})>S_0+a$ or $j(\alpha)=j(\beta)$, then $w_{\alpha\beta}=0$. Otherwise let $\mathcal E$ be the closed ``ellipse''
\[
d(\mathring U_{j(\alpha)},x)+d(\mathring U_{j(\beta)},x)\le S_0+a.
\]
Choose an open set $\Omega$ with smooth boundary such that $U_{j(\alpha)}\subset\Omega$, $U_{j(\beta)}\cap\Omega=\varnothing$, with the two pieces of $\mathcal E$ lying in the corresponding $M_j$'s, and put $\Gamma=\partial\Omega\cap\mathcal E$. Then for every $\eps>0$,
\begin{equation}
w_{\alpha\beta}=h^2\int_\Gamma\left(\varphi_\alpha\frac{\partial\varphi_\beta}{\partial n}
-\varphi_\beta\frac{\partial\varphi_\alpha}{\partial n}\right)dS
\quad\bmod\ \cO\!\left(e^{-(S_0+a-\eps)/h}\right),
\end{equation}
where $n$ is the interior normal with respect to $\Omega$.
\end{remark}
Indeed, one inserts a cutoff $\chi_{\mathcal E}$ equal to $1$ on $\mathcal E$ and uses Stokes' formula
\[
\int_\Omega\omega\cdot\nabla u\,dx
=\int_\Omega(\nabla^*\omega)u\,dx-
\int_{\partial\Omega}\left\langle\omega,\frac\partial{\partial n}\right\rangle u\,dS,
\]
together with
\[
\nabla^*(\varphi_\alpha\nabla\varphi_\beta-
\varphi_\beta\nabla\varphi_\alpha)
=-\varphi_\alpha\Delta\varphi_\beta+
\varphi_\beta\Delta\varphi_\alpha=0.
\]

The same arguments can be applied with rapidly decreasing errors. Assume, for every $N\in\N$,
\begin{equation}
\mu_\alpha-\mu_\beta=\cO(h^N),
\tag{2.24'}
\end{equation}
and assume the improved energy decay statement
\begin{equation}
\begin{minipage}{0.84\textwidth}
For every $\alpha$ and every compact $K\subset\mathring M_{j(\alpha)}$, there is $N_0$ such that
\[
\norm{\nabla(e^{\phi_{j(\alpha)}/h}\varphi_\alpha)}_{L^2(K)}+
\norm{e^{\phi_{j(\alpha)}/h}\varphi_\alpha}_{L^2(K)}=\cO(h^{-N_0}).
\]
\end{minipage}
\end{equation}
Then, if $d(\mathring U_{j(\alpha)},\mathring U_{j(\beta)})>S_0$ or $\alpha=\beta$, one has
\[
w_{\alpha\beta}=\cO(h^Ne^{-S_0/h})\qquad\text{for every }N.
\]
Otherwise, choosing $a>0$ small and $\Gamma$ as above,
\begin{equation}
w_{\alpha\beta}=h^2\int_\Gamma\left(\varphi_\alpha\frac{\partial\varphi_\beta}{\partial n}
-\varphi_\beta\frac{\partial\varphi_\alpha}{\partial n}\right)dS
+\cO(h^Ne^{-S_0/h}),\qquad N\in\N.
\tag{2.25'}
\end{equation}
In this situation $w_{\alpha\beta}=w_{\beta\alpha}=\widehat w_{\alpha\beta}$ modulo $\cO(h^Ne^{-S_0/h})$ for every $N$.

\begin{corollary}
In the special case when there is only one well $U_j$, then for $h$ sufficiently small there is a bijection
\[
b:\Sp(P)\cap I(h)\longrightarrow\Sp(P_{M_j})\cap I(h)
\]
such that $|b(\lambda)-\lambda|=\cO(e^{-2\sigma/h})$ for every $\sigma<S_1$.
\end{corollary}
The same discussion is valid for a Dirichlet problem when $M$ is a compact Riemannian manifold with boundary and $\mathring M_j\subset M$.

Returning to the general situation, replace $M_1$ by another admissible domain $\widetilde M_1\Subset M$ containing $B(U_1,S_1)$ and disjoint from the other wells, where
\[
S_1<S^{(1)}:=\min_{k\ne1}d(U_1,U_k).
\]

\begin{proposition}
Let $e(h)\in(0,a(h))$ and $\log e(h)=o(1/h)$. Then for $h$ sufficiently small there is a bijection
\[
b:\Sp(P_{M_1})\cap I(h)\longrightarrow
\Sp(P_{\widetilde M_1})\cap(I(h)+B(0,e(h)))
\]
such that
\[
|b(\lambda)-\lambda|=\cO(e^{-2\sigma/h})\qquad\text{for every }\sigma<S_1.
\]
\end{proposition}
\begin{proof}
Introducing a third domain $M'_1$ reduces the proof to the case $M_1\Subset M'_1$. Let $J(h)$ have the same spectral separation properties as $I(h)$ for $P$, $P_{M_1}$ and $P_{M'_1}$ and satisfy $I(h)+B(0,2e(h))\subset J(h)$. Corollary~2.14 gives a bijection between the two Dirichlet spectra in $J(h)$, and it suffices to restrict it to $I(h)\cap\Sp(P_{M_1})$.
\end{proof}

We now define the ``spectrum'' $\Sp(U_1)$ as the collection of the sets $\Sp(P_{M_1})\cap(I(h)+B(0,c(h)))$ occurring in the preceding proposition. If $h\mapsto\lambda(h)\in\R$ and $\sigma<S^{(1)}$, we write
\[
d(\lambda,\Sp(U_1))=\Ot(e^{-2\sigma/h})
\]
when the distance from $\lambda(h)$ to the corresponding Dirichlet spectrum is $\cO(e^{-2\sigma/h})$ for every admissible choice of $M_1$ containing $B(U_1,\sigma)$. The same terminology will be used for the other wells.

Let $u(h)\in D(P)$, $\norm u=1$, $(P-\lambda)u=0$, with $\lambda(h)\in I(h)$. We consider functions $\psi$ on $M$ such that
\begin{equation}
|\psi(x)-\psi(y)|\le d(x,y),\qquad x,y\in M,
\end{equation}
with $\psi$ constant on each $U_j$, and such that
\begin{equation}
\norm{\nabla(e^{\psi/h}u)}_{L^2(K)}+
\norm{e^{\psi/h}u}_{L^2(K)}=\cO(e^{\eps/h})
\end{equation}
for every $\eps>0$ and compact $K\subset M$. If $\psi_1,\psi_2$ have these properties, then so does $\max(\psi_1,\psi_2)$. We assume that the $U_j$ have diameter zero for the Agmon metric and that
\begin{equation}
\begin{minipage}{0.84\textwidth}
$\psi$ is the largest function satisfying (2.27), (2.28), and no larger function can be found even if (2.28) is required only for $\eps$ in a subset having $0$ as an accumulation point.
\end{minipage}
\end{equation}
Put $a_j=\psi|_{U_j}$. At least one $a_j$ is zero and $|a_j-a_k|\le d(U_j,U_k)$. As in Lemmas~2.3 and~2.7,
\begin{equation}
\psi(x)=\min_j\bigl(d(x,U_j)+a_j\bigr).
\end{equation}
For each $j$, let $r_j$ be the largest number in $[0,S_j]$ such that $\psi(x)=a_j+d(x,U_j)$ on $B(U_j,r_j)$. If $a_j=0$, then $r_j\ge S_0/2$. If $a_j>0$, comparison on $\partial B(U_j,r_j)$ gives
\[
a_j+2r_j\le a_k+d(U_j,U_k),\qquad k\ne j.
\]
Moreover, for every $\eps>0$ there is $C_\eps$ such that
\begin{equation}
C_\eps^{-1}e^{-\eps/h}\le
e^{a_j/h}\left(\norm{\nabla u}_{L^2(B(U_j,r_j))}+
\norm u_{L^2(B(U_j,r_j))}\right)\le C_\eps e^{\eps/h}.
\end{equation}
The right inequality follows from (2.28), while failure of the left inequality would allow one to improve $\psi$, contradicting (2.29). Standard a priori estimates allow the gradient term to be removed.

\begin{proposition}
Under the above assumptions, for every $\eps>0$,
\[
d(\lambda,\Sp(U_j))=\cO_\eps\!\left(e^{-2(r_j-\eps)/h}\right).
\]
\end{proposition}
\begin{proof}
Choose $\chi\in C_0^\infty(B(U_j,r_j))$ equal to $1$ near $B(U_j,r_j-\eps)$. Let $\widetilde u$ be a normalization of $u$ in $L^2(B(U_j,r_j))$, and set $v=\chi\widetilde u$. Then $\norm v=1$ and
\begin{equation}
P_jv=\lambda v+\cO\!\left(e^{-(r_j-\eps)/h}\right)
\quad\text{in }L^2(B(U_j,r_j)).
\end{equation}
Let $F\subset L^2(M_j)$ be the sum of the eigenspaces associated with $\Sp(P_{M_j})\cap I(h)$. As in Proposition~2.5,
\[
d(v,F)=\cO_\eps(e^{-(r_j-\eps)/h}),
\]
so $\norm{\Pi_Fv}=1$ modulo exponentially small errors. Since
\[
P\Pi_Fv=\lambda\Pi_Fv+\Pi_F[P,\chi]u,
\]
and expansion in an orthonormal eigenbasis of $F$ together with the decay estimates gives
\[
\norm{\Pi_F[P,\chi]u}=\cO_\eps(e^{-2(r_j-\eps)/h}),
\]
the claim follows.
\end{proof}

If $a_j=0$, then the source records
\[
r_j\ge\frac12\min_{k\ne j}d(U_j,U_k)=\frac14S_j,
\]
and hence
\[
d(\lambda,\Sp(U_j))=\cO_\eps(e^{-(S_j-\eps)/h}).
\]
On the other hand, if for some $j$
\[
d(\lambda,\Sp(U_j))>C_\eps^{-1}e^{-\eps/h}
\]
for every $\eps>0$, then $r_j=0$. We then say that $U_j$ is a \emph{non-resonant well}. In a future paper the authors plan to study the interaction of a group of wells ``through'' a second group of non-resonant wells.

\section{Formal constructions near a strict minimum of the potential}
The computations of this section are in principle well known. They correspond to the symplectic case in the study of operators with double characteristics. B. Simon [22] has very recently given results in the present framework, showing that the lowest eigenvalues in the case of one point-well, where $V$ has a strict minimum, admit asymptotic expansions in powers of $h$. It is not true, however, that only integer powers of $h$ appear in these expansions. We shall also need more precise information about the eigenfunctions and, in this section and the following, we follow our original plan to use a canonical transformation. This transformation is not absolutely necessary, but it makes the computations simpler and permits us to examine explicitly when half exponents appear. In the next section we prove that in the analytic case the asymptotic expansions are analytic symbols, and in Section~5 we combine these results with decay estimates to link the formal eigenvalues and eigenfunctions with the actual ones.

The general assumption in this and the following two sections is that one $M_j$ is fixed and the corresponding well is reduced to a point, $U_j=\{x_0\}$. The potential has a strict minimum at $x_0$:
\begin{equation}
V''(x_0)>0.
\end{equation}
To $P$ we associate the symbol
\[
p(x,\xi)=\xi^2+V(x),
\]
where $\xi^2$ denotes $\norm\xi^2$ (extended holomorphically in the Riemannian case). The weights $e^{\psi(x)/h}$ correspond to the complex phases $i\psi(x)$, and it is therefore natural to consider
\[
-p(x,i\xi)=\xi^2-V(x)=q(x,\xi).
\]
Let $\lambda_1,\ldots,\lambda_n$ be the eigenvalues of the quadratic form $\frac12\langle V''(x_0)t,t\rangle$ with respect to the Riemannian metric. Coordinates centered at $x_0$ can be chosen so that
\begin{equation}
q(x,\xi)=\xi^2-\sum_{j=1}^n\lambda_jx_j^2+\cO(|(x,\xi)|^3),
\end{equation}
and
\begin{equation}
H_q=2\left(\xi\cdot\frac\partial{\partial x}+
\sum_{j=1}^n\lambda_jx_j\frac\partial{\partial\xi_j}\right)
+\cO(|(x,\xi)|^2).
\end{equation}
The linearization at $(0,0)$ is the fundamental matrix
\begin{equation}
F=2\begin{pmatrix}0&I\\ \Lambda&0\end{pmatrix},
\qquad \Lambda=\operatorname{diag}(\lambda_1,\ldots,\lambda_n).
\end{equation}
Its eigenvalues are $\pm2\sqrt{\lambda_j}$ and the corresponding stable and unstable eigenspaces are
\begin{equation}
\Lambda_\pm^0:\qquad \xi_j=\pm\sqrt{\lambda_j}\,x_j.
\end{equation}
The spaces $\Lambda_\pm^0$ are Lagrangian and have the form
\begin{equation}
\xi=\pm\phi'(x),
\end{equation}
where, in the chosen coordinates,
\[
\phi(x)=\frac12\sum_{j=1}^n\sqrt{\lambda_j}\,x_j^2.
\]

By the stable manifold theorem there are, in a suitable neighborhood $W$ of $(0,0)$, unique closed connected $C^\infty$ manifolds $\Lambda_\pm$ of dimension $n$ such that
\begin{equation}
T_{(0,0)}\Lambda_\pm=\Lambda_\pm^0,
\qquad H_q\text{ is tangent to }\Lambda_\pm.
\end{equation}
Moreover
\begin{equation}
\Lambda_\pm=\{(x,\xi)\in W;\ \exp(-tH_q)(x,\xi)\in W\text{ for all }t\ge0\},
\end{equation}
and points of $\Lambda_\pm$ tend exponentially to $(0,0)$ under the corresponding flow. The same holds in the real analytic category. The manifolds are Lagrangian and hence
\begin{equation}
\Lambda_\pm:\qquad \xi=\pm\phi'(x),\qquad x\in\Omega,
\end{equation}
where $\phi(0)=0$ and $\phi'(0)=0$.

Let $t\mapsto(x(t),\xi(t))\in\Lambda_+$ be an integral curve of $H_q$. Then
\[
\frac d{dt}\phi(x(t))=\phi'(x(t))\cdot\dot x(t)
=\xi(t)\cdot\dot x(t)=2\xi(t)\cdot\xi(t)
=\sqrt{V(x(t))}\,\norm{\dot x(t)}.
\]
Hence the length of the projected curve for the Agmon metric is
\begin{equation}
\int_a^b\sqrt{V(x(t))}\,\norm{\dot x(t)}\,dt
=\phi(x(b))-\phi(x(a)).
\end{equation}
After restricting $\Omega$, write $\Omega_y=\{x\in\Omega;\ \phi(x)<y\}$ with $\Omega_y\Subset\Omega_{y'}$ for $0<y<y'<y_0$.

\begin{proposition}
Under the assumptions above,
\[
\phi(x)=d(x,0),\qquad x\in\Omega.
\]
\end{proposition}
\begin{proof}
Let $x\in\Omega$. To prove $d(x,0)\ge\phi(x)$, take a $C^1$ curve $\gamma:[a,b]\to M$ with $\gamma(a)=0$ and $\gamma(b)=x$. Restricting to a segment if necessary, suppose $\phi(\gamma(b))=y$. Then
\[
y=\int_a^b\phi'(\gamma(t))\cdot\dot\gamma(t)\,dt
\le\int_a^b\norm{\phi'(\gamma(t))}\,\norm{\dot\gamma(t)}\,dt
=\operatorname{length}(\gamma),
\]
since $\norm{\phi'}^2=V$. Hence $\phi(x)\le d(x,0)$. Conversely, the projection of the bicharacteristic in $\Lambda_+$ from $(0,0)$ to $(x,\phi'(x))$ has length $\phi(x)$; approximating its infinite endpoint at the origin by finite curves gives $d(x,0)\le\phi(x)$.
\end{proof}

We now choose local coordinates by the Morse lemma so that
\[
\phi(x)=\frac12x^2
\]
and $V''(0)$ is diagonal. Introduce the complex Lagrangian manifolds $L_\pm$ by $\xi=\pm i\phi'(x)=\pm ix$. Consider the phase function
\[
f(x,y)=\frac i2(x-y)^2-\frac14x^2
\]
and the associated complex canonical transformation
\[
\mathcal H_f:(y,i(x-y))\longmapsto
\left(x,i(x-y)-\frac i2x\right).
\]
Then $\mathcal H_f(L_+)$ is given by $\xi=0$, $\mathcal H_f(L_-)$ by $x=0$, and $\mathcal H_f(\R^{2n})$ is the $I$-Lagrangian manifold
\[
\xi=\frac2i\frac{\partial\Phi}{\partial x},
\qquad \Phi(x)=\frac14|x|^2.
\]
We define the FBI transform, with $h=1/\lambda$, by
\begin{equation}
T_1u(x,\lambda)=\lambda^{n/2}C_n\int e^{i\lambda f(x,y)}u(y)\,dy.
\end{equation}
It maps $\mathcal S'(\R^n)$ into the local weighted holomorphic space $H^{\mathrm{loc}}_{\Phi,0}(\C^n)$: for each compact $K$ and each $\eps>0$,
\[
\sup_K|v(x,\lambda)|e^{-\lambda\Phi(x)}=\cO(e^{\eps\lambda}).
\]
The constant $C_n>0$ is chosen so that $T_1$ is formally unitary. If $u(y,\lambda)$ is an analytic symbol near $0$ and $\widetilde u(y,\lambda)=u(y,\lambda)e^{-\lambda\phi(y)}$, then $T_1\widetilde u$ is well defined, modulo an exponentially decreasing term, by restricting the integration to a neighborhood of $0$. The map
\[
H_{0,0}\ni u\longmapsto T_1(ue^{-\lambda\phi})\in H^{\mathrm{loc}}_{\Phi,0}
\]
is a bijection modulo exponentially small terms, and stationary phase preserves classical analytic symbols.

Recall the symbol terminology. A formal classical analytic symbol of order $m$ in $\Omega\subset\C^n$ is an expression
\[
\sum_{j\in\N/2}\lambda^{m-j}u_j(x),
\]
where the $u_j$ are holomorphic and, for every $K\Subset\Omega$, $|u_j(x)|\le C^{j+1}j!$. A realization is obtained by an asymptotic sum over $j\le1/(\eps C)$. We also use half-degree expansions indexed by $\frac12\N$.

For $m\in\R$, $k\in\N$, let $S^{m,k}$ and $S_{1/2}^{m,k}$ be the spaces of formal symbols
\[
\sum_j u_j(x)\lambda^{m-j},
\]
where $u_j$ is a formal power series satisfying $u_j=\cO(|x|^{k-2j})$, with $j\in\N$ or $j\in\frac12\N$, respectively. We have
\[
S^{m,k}\subset S^{m',k'}
\]
when $m'-m\in\N$ (or $\frac12\N$ in the half-order case) and $m-k/2\le m'-k'/2$. The total order is at most $m-k/2$. Put
\[
S=\bigcup_{m\in\frac12\N}S^{m,0},
\qquad
S_{1/2}=\bigcup_{m\in\frac12\N}S_{1/2}^{m,0}.
\]
The FBI transform sends these symbol classes into the corresponding classes on the transform side.

Let $S^m(\R_+)$ denote the formal classical symbols
\[
\sum_{j\in\N/2}a_j\lambda^{m-j},\qquad a_j\in\C,
\]
and put
\[
S_{1/2}^m(\R_+)=S^m(\R_+)+S^{m-1/2}(\R_+).
\]
If $v_j\in S^{m_j,k_j}$, define the formal weighted scalar product
\begin{equation}
(v_1\mid v_2)_\Phi=\lambda^n\int
v_1(x,\lambda)\overline{v_2(x,\lambda)}e^{-2\lambda\Phi(x)}L(dx).
\end{equation}
By formal stationary phase it belongs to
\[
S^{m_1+m_2-(k_1+k_2)/2}
\]
if $k_1+k_2$ is even, and to
\[
S^{m_1+m_2-(k_1+k_2+1)/2}
\]
if the sum is odd; analogous statements hold in the half-degree case.

If $C_n$ is suitably chosen, then for all $u_1,u_2$ in the symbol classes,
\begin{equation}
\bigl(T_1(u_1e^{-\lambda\phi})\mid T_1(u_2e^{-\lambda\phi})\bigr)_\Phi
=\lambda^{n/2}\int u_1(y,\lambda)\overline{u_2(y,\lambda)}e^{-2\lambda\phi(y)}\,dy.
\end{equation}
For completeness, the formal verification is as follows. One has
\[
\Phi(x)=\sup_{t\in\R^n}(-\operatorname{Im}f(x,t))
=\frac12\,\vc_t\bigl(f(x,t)-\overline{f(x,\bar t)}\bigr).
\]
Put $g(x,t)=-\overline{f(x,\bar t)}$ and define
\begin{equation}
\frac2i\psi(x,y)=\vc_t\bigl(f(x,t)+g(y,t)\bigr).
\end{equation}
The function $\psi$ is the unique entire function with $\psi(x,\bar x)=\Phi(x)$. Since $\det \partial_x\partial_t f\ne0$, the relation can be inverted:
\begin{equation}
-g(y,t)=\vc_x\left(f(x,t)-\frac2i\psi(x,y)\right).
\tag{3.14'}
\end{equation}
\addtocounter{equation}{-1}
With $\widetilde u_j=u_je^{-\lambda\phi}$, stationary phase eliminates the $z$-integration in the scalar product, and the Kuranishi trick reduces the remaining $(y,w)$-integration to a constant multiple of the identity. Choosing $C_n$ to normalize this constant gives (3.13).

If the Riemannian volume form is $a(y)dy$, $a>0$, put
\begin{equation}
Tu=T_1(a^{1/2}u).
\end{equation}
Then
\begin{equation}
\bigl(T(u_1e^{-\lambda\phi})\mid T(u_2e^{-\lambda\phi})\bigr)_\Phi
=\lambda^{n/2}\int u_1\overline{u_2}e^{-2\lambda\phi}a\,dy.
\tag{3.13'}
\end{equation}
\addtocounter{equation}{-1}

In the analytic case, [24] gives a formal classical analytic symbol
\[
\widetilde P(x,\xi,\lambda)\sim\sum_j\widetilde p_j(x,\xi)\lambda^{-j}
\]
near $(0,0)$ such that, after choosing a realization and the corresponding pseudodifferential operator,
\begin{equation}
\widetilde P\,T=T\,P
\end{equation}
modulo exponentially small terms. The principal symbol satisfies
\begin{equation}
\widetilde p\circ\mathcal H_f=p.
\end{equation}
Since $p$ vanishes on $L_+$, $\widetilde p$ vanishes for $x=0$ and for $\xi=0$, hence
\begin{equation}
\widetilde p(x,\xi)=i\sum_{|\alpha|=|\beta|=1}
x^\alpha a_{\alpha\beta}(x,\xi)\xi^\beta.
\end{equation}
By symplectic invariance of the fundamental matrix,
\begin{equation}
\sum_{|\alpha|=|\beta|=1}x^\alpha a_{\alpha\beta}(0,0)\xi^\beta
=\sum_{j=1}^n\mu_jx_j\xi_j,
\qquad \mu_j=2\sqrt{\lambda_j}.
\end{equation}
The transformed operator acts on the symbol classes and
\begin{equation}
\widetilde P(Tu)=T(Pu).
\end{equation}
In the $C^\infty$ case the coefficients are only formal Taylor series, but the same relations remain valid. Formal self-adjointness follows from that of $P$ and the scalar product above. The subprincipal symbol satisfies
\begin{equation}
\widetilde p_{-1}(0,0)=\frac12\sum_{j=1}^n\mu_j.
\end{equation}

We now study $Q=\lambda^{-1}\widetilde P$ on $S^{m,k}$ and $S_{1/2}^{m,k}$. We can write formally
\[
Q=\sum_{|\alpha|=|\beta|=1}x^\alpha a_{\alpha\beta}(x,D_x)x^\beta+a_0(x,D_x),
\]
and put
\[
Q_0=\sum_{|\alpha|=|\beta|=1}x^\alpha a_{\alpha\beta}(0,0)x^\beta+a_0(0,0).
\]
The operator $Q_0$ preserves the homogeneous spaces. If $S_h^{m,k}$ denotes the space of elements $\lambda^m u_m(x)$ with $u_m$ homogeneous of degree $k$, then $Q_0:S_h^{m,k}\to S_h^{m,k}$, while
\[
Q-Q_0:S^{m,k}\to S^{m,k+1}
\]
and similarly in the half-degree classes.

The monomials are eigenfunctions of $Q_0$; the eigenvalue corresponding to $x^\alpha$ is
\[
\sum_j\left(\alpha_j+\frac12\right)\mu_j.
\]
Let $E_0$ be such an eigenvalue and choose a basis $e_1,\ldots,e_{N_0}$ of eigenvectors with $e_j\in S_h^{m_j,k_j}$ and $m_j-k_j/2=0$. If $u\in S_{1/2}^{m,k}$ and $(Q_0-E_0)u=0$, then
\[
u=\sum_{j=1}^{N_0}\lambda_j e_j
\]
with uniquely determined scalar symbols $\lambda_j$. If
\[
E\notin\left\{\sum_j\left(\alpha_j+\frac12\right)\mu_j;\ \alpha\in\N^n\right\}=\Sp(Q_0),
\]
then $Q_0-E$ is invertible on the polynomial and symbol spaces.

Let $\gamma$ be a simple loop containing $E_0$ and no other point of $\Sp(Q_0)$. Define
\begin{equation}
\Pi_0=\frac1{2\pi i}\int_\gamma(E-Q_0)^{-1}\,dE.
\end{equation}
This projects onto $\ker(Q_0-E_0)$ in each homogeneous class. Similarly,
\begin{equation}
\Pi=\frac1{2\pi i}\int_\gamma(E-Q)^{-1}\,dE.
\end{equation}
The resolvent is given asymptotically by the Neumann series
\[
(Q-E)^{-1}=(Q_0-E)^{-1}+(Q_0-E)^{-1}(Q_0-Q)(Q_0-E)^{-1}+\cdots,
\]
so $\Pi-\Pi_0$ raises the regularizing degree by $1/2$, and $[\Pi,Q]=0$.

\begin{lemma}
If $F\in\Sp(Q_0)$ lies outside $\gamma$, then
\begin{equation}
(F-Q)^{-1}\Pi=\frac1{2\pi i}\int_\gamma
(F-E)^{-1}(E-Q)^{-1}\,dE.
\end{equation}
\end{lemma}
\begin{proof}
Let $A$ denote the right hand side. Multiplication by $F-Q$ gives
\[
(F-Q)A=\frac1{2\pi i}\int_\gamma(F-E)^{-1}\,dE\,I+\Pi=\Pi,
\]
so $A=(F-Q)^{-1}\Pi$.
\end{proof}

\begin{lemma}
$\Pi$ is a projection: $\Pi^2=\Pi$.
\end{lemma}
\begin{proof}
Let $\Gamma$ be a loop with the same spectral separation properties as $\gamma$ and containing $\gamma$. By Lemma~3.2,
\[
\Pi^2=\frac1{2\pi i}\int_\Gamma(F-Q)^{-1}\Pi\,dF
=\frac1{(2\pi i)^2}\int_\Gamma\int_\gamma
(F-E)^{-1}(E-Q)^{-1}\,dE\,dF=\Pi.
\]
\end{proof}

\begin{proposition}
The vectors $\Pi(e_1),\ldots,\Pi(e_{N_0})$ form a basis of $\Pi(S_{1/2})$ in the following sense.
\begin{enumerate}[label=(\alph*)]
\item If $u\in\Pi(S_{1/2})$ has total degree at most $m_0$, then
\[
u=\sum_{j=1}^{N_0}\lambda_j\Pi(e_j),
\qquad \lambda_j\in S_{1/2}^{m_0}(\R_+).
\]
\item If $\sum_j\lambda_j(\Pi(e_j)+v_j)=0$ with $v_j$ of lower order, then all $\lambda_j$ vanish.
\end{enumerate}
\end{proposition}
\begin{proof}
For (a), choose leading coefficients $\lambda_j^0$ so that $\Pi_0u=\sum\lambda_j^0e_j$. Then
\[
\Pi\Pi_0u=\sum_j\lambda_j^0\Pi(e_j)
=u+\Pi(\Pi_0-\Pi)u=u-u^{(1)},
\]
where $u^{(1)}$ has total degree lower by $1/2$. Iterating gives
\[
u^{(1)}=\sum_j\lambda_j^1\Pi(e_j)+u^{(2)},\qquad
\lambda_j^1\in S_{1/2}^{m_0-1/2},
\]
and hence $\lambda_j\sim\lambda_j^0+\lambda_j^1+\cdots$. Part (b) follows by taking the highest homogeneous part and iterating downward.
\end{proof}

After rearranging the $\Pi(e_j)$ we may assume that
\begin{equation}
k_j\text{ is even for }1\le j\le N_0,
\qquad k_j\text{ is odd for }N_0+1\le j\le\widetilde N_0.
\end{equation}
The Gram matrix of the $\Pi(e_j)$ has the form $I+A$, where, with respect to this parity decomposition,
\[
A=\begin{pmatrix}A_{11}&A_{12}\\A_{21}&A_{22}\end{pmatrix},
\qquad
A_{11},A_{22}\in S^{-1},\quad A_{12},A_{21}\in S^{-1/2}.
\]
By a binomial series construct $(I+A)^{-1/2}=I+B$, with $B=B^*$ of the same block structure, and define
\[
\begin{pmatrix}f_1\\ \vdots\\ f_{N_0}\end{pmatrix}
=(I+B)\begin{pmatrix}\Pi(e_1)\\ \vdots\\ \Pi(e_{N_0})\end{pmatrix}.
\]
Then $(f_j)$ is an orthonormal basis of $\Pi(S_{1/2})$, still satisfying Proposition~3.4.

\begin{proposition}
Each $f_j$ belongs to a class $S_{1/2}^{\widetilde m_j,\widetilde k_j}$ with $\widetilde m_j-\widetilde k_j/2=0$. The integer $\widetilde k_j$ is even in the first parity group and odd in the second.
\end{proposition}
\begin{proof}
Each $f_j$ is a sum of terms of total degree zero. Terms coming from a half-order off-diagonal block can be reindexed by decreasing both the symbol order and polynomial degree by $1/2$ and $1$, respectively, without changing the total order. Thus all terms can be placed in a common symbol class of the asserted parity.
\end{proof}

We have $Qf_k\in S^{\widetilde m_k,\widetilde k_k}$ and $Qf_k=E_0f_k$ modulo one higher regularizing degree. Hence the matrix of $Q$ in the orthonormal basis $f_1,\ldots,f_{N_0}$ has the form
\begin{equation}
((Qf_k\mid f_j))=E_0I+
\begin{pmatrix}Q_{11}&Q_{12}\\Q_{21}&Q_{22}\end{pmatrix},
\end{equation}
where $Q$ is Hermitian, $Q_{11},Q_{22}\in S^{-1}(\R_+)$ and $Q_{12},Q_{21}\in S^{-1/2}(\R_+)$.

\begin{theorem}
There exists, in the sense of Proposition~3.4, an orthonormal basis of eigenvectors of $Q$ in $\Pi(S_{1/2})$,
\[
g_j\in S_{1/2}^{\widehat m_j,\widehat k_j},
\qquad \widehat m_j=\frac12\widehat k_j,
\]
and the corresponding eigenvalues belong to $E_0+S_{1/2}^{-1/2}(\R_+)$. If only one parity occurs in (3.25), the $g_j$ may be taken in integer-degree symbol classes and the eigenvalues belong to $E_0+S^{-1}(\R_+)$. The corresponding $\widehat k_j$ have the parity of the surviving group.
\end{theorem}
\begin{proof}
Subtract $E_0$ and write the finite Hermitian matrix as
\[
Q=\sum_{\nu\ge1/2,\ \nu\in\frac12\N}h^\nu Q_\nu.
\]
Its truncations
\[
Q_N=\sum_{1/2\le\nu\le N}h^\nu Q_\nu
\]
are analytic functions of $h^{1/2}$, so their eigenvalues are analytic functions of $h^{1/2}$. After consistent ordering, the eigenvalues of higher truncations differ by $\cO(h^N)$, giving formal asymptotic eigenvalues in $S_{1/2}^{-1/2}$; they are the roots of $\det(Q-E)$.

For a formal eigenvalue of multiplicity $d$, choose $N_1$ large enough that for $N>N_1$ the associated $d$ eigenvalues form a cluster of diameter $\cO(h^N)$ separated from the rest by a power of $h$. The corresponding spectral projection is a meromorphic function of $h^{1/2}$, bounded for real positive $h$, and hence belongs to $S_{1/2}^0$. Starting with an orthonormal basis at level $N_1$, project it to the higher truncations and orthonormalize. The resulting sequences define formal eigenvectors $g_j$ forming an orthonormal basis of $\ker(Q-\lambda_j)$. The integer-degree case is identical, starting from the improved block order $S^{-1}$.
\end{proof}

We finally return to $P$.

\begin{theorem}
Let $E_0$ be as above. There exist amplitudes
\[
a_j\in S_{1/2}^{\widehat m_j,\widehat k_j},
\qquad \widehat m_j=\frac12\widehat k_j,
\]
such that $u_j=a_je^{-\lambda\phi}$ form an orthonormal system for the formal scalar product corresponding to (3.15$'$), and
\[
Pu_j=E_ju_j,
\qquad E_j\in E_0+S_{1/2}^{-1/2}(\R_+).
\]
If only one parity occurs, the half-degree terms disappear as in Theorem~3.6.
\end{theorem}

\begin{remark}
One can avoid FBI transforms in the proof by working directly with $Q=\lambda e^{\lambda\phi}Pe^{-\lambda\phi}$ in suitable coordinates. The leading operator then has a less pleasant form and its eigenfunctions are Hermite polynomials rather than homogeneous polynomials; see Damburg--Propin [28].
\end{remark}

\begin{remark}
A sufficient condition for the disappearance of half-degree terms in the eigenvalues is that, for all $\alpha\in\N^n$ with
\[
\sum_j\left(\alpha_j+\frac12\right)\mu_j=E_0,
\]
the integers $\sum_j\alpha_j$ have the same parity.
\end{remark}

\begin{remark}
Write formally
\[
a_j\sim\sum_{\nu\ge0}a_{j,\widehat m_j-\nu}(x)\lambda^{\widehat m_j-\nu}
\]
and choose smooth functions $\widehat a_{j,k}$ having the corresponding Taylor series at $0$. Then the formal classical symbol
\[
\widehat a_j(x,\lambda)\sim\sum_{\nu\ge0}\widehat a_{j,\widehat m_j-\nu}(x)\lambda^{\widehat m_j-\nu}
\]
can be modified by terms flat at $0$ so that the transport equations are solved exactly in a sufficiently small neighborhood. The correction reduces to repeated equations
\[
(\nu(x,\partial_x)-f(x))u(x)=v(x),
\]
where $v=\cO(|x|^N)$ and $\nu$ is a real vector field vanishing only at $0$, whose linearization has strictly positive eigenvalues. If $x_t=\exp(t\nu)(x)$, then
\[
\norm{D_x^\alpha x_t}\le C_\alpha e^{-|t|/C}|x|,
\qquad t\le0,
\]
and the flat solution is
\begin{equation}
u(x)=\int_{-\infty}^0v(x_t)
\exp\left(\int_t^0 f(x_s)\,ds\right)dt.
\end{equation}
Thus one can find formal classical symbols $\widehat a_j$ representing the amplitudes of Theorem~3.7 so that
\[
(\lambda P-\widehat E_j)(\widehat a_je^{-\lambda\phi})=0
\]
in a small neighborhood. If $\widetilde a_j$ and $\widetilde E_j$ are asymptotic sums, then
\begin{equation}
(\lambda P-\widetilde E_j)(\widetilde a_je^{-\lambda\phi})
=\widetilde b(x,\lambda)e^{-\lambda\phi},
\qquad \widetilde b\in S^{-\infty}(\Omega\times\R_+).
\end{equation}
\end{remark}

\begin{remark}
Complex conjugation commutes with the FBI transform, with $P$, and hence with the transformed operator. The monomial basis may be chosen real; the spectral projection, Gram matrices, orthonormalization, eigenvectors, amplitudes, and their smooth realizations can therefore all be chosen real valued.
\end{remark}

We end the section with an example in which half degrees genuinely occur. On the FBI-transform side consider
\begin{equation}
Q=2x_1\partial_{x_1}+x_2\partial_{x_2}
+x_2(a x_2+\bar a\,\partial_{x_2})\partial_{x_1}
+x_1(a x_2+\bar a\,\partial_{x_2})\partial_{x_2}.
\tag{3.28}
\end{equation}
Assume coordinates have been chosen so that $\phi=\frac12|x|^2$. The Weyl symbol of $Q$ restricted to $\Lambda_\phi$ is real, so $Q$ is self-adjoint for the Bargmann--Fock scalar product. After $x=\lambda^{-1/2}y$,
\[
Q=Q_0+\lambda^{-1/2}Q_1,
\]
where
\[
Q_0=2y_1\partial_{y_1}+y_2\partial_{y_2},
\qquad
Q_1=y_2(a y_2+\bar a\,\partial_{y_2})\partial_{y_1}
+y_1(a y_2+\bar a\,\partial_{y_2})\partial_{y_2}.
\]
A symbol $u=u_0+\lambda^{-1/2}u_1+\lambda^{-1}u_2+\cdots$ is an eigenvector with $E=E_0+\lambda^{-1/2}E_1+\lambda^{-1}E_2+\cdots$ precisely when
\begin{align*}
(Q_0-E_0)u_0&=0,\\
(Q_0-E_0)u_1+(Q_1-E_1)u_0&=0,\\
(Q_0-E_0)u_2+(Q_1-E_1)u_1&=0,\quad\ldots
\end{align*}
In the Bargmann--Fock space with scalar product
\begin{equation}
(u\mid v)=\int u(x)\overline{v(x)}e^{-2\phi(x)}L(dx),
\end{equation}
the monomials form an orthogonal basis. If $\mathcal E_0$ is the eigenspace for $E_0$ and $\Pi_0$ the orthogonal projection onto it, the first correction equation is solvable if and only if
\[
(\Pi_0Q_1-E_1)u_0=0.
\]
Thus if $\Pi_0Q_1|_{\mathcal E_0}$ is non-degenerate, then $E_1\ne0$ and half degrees appear.

For $E_0=2$, $\mathcal E_0$ is spanned by $y_1$ and $y_2^2$. One computes
\[
Q_1(y_1)=a y_2^3,
\qquad
Q_1(y_2^2)=2a y_2^2y_1+2\bar a\,y_1,
\]
so in the basis $(y_1,y_2^2)$,
\[
\Pi_0Q_1=\begin{pmatrix}0&2\bar a\\ a&0\end{pmatrix},
\]
which is non-degenerate if $a\ne0$.

The same phenomenon occurs directly for the Schr\"odinger operator
\[
h^2D_{x_1}^2+x_1^2+2(h^2D_{x_2}^2+x_2^2)+V(x),
\]
where $V$ is homogeneous of degree $3$. Under $x_j=h^{1/2}y_j$ this becomes
\[
h\left(D_{y_1}^2+y_1^2+2(D_{y_2}^2+y_2^2)+h^{1/2}V(y_1,y_2)\right).
\]
The unperturbed operator has a double eigenvalue, and $E_1\ne0$ for example when $V(x)=x_1^2x_2$.

\section{Further information in the analytic case}
In this section we assume some familiarity with analytic pseudodifferential operators and use the terminology and results of [24]. Throughout the section $M$ and $V$ are real analytic and the other assumptions of Section~3 remain in force. The distance function $\phi(x)$ is analytic near $0$, and the amplitude in the FBI transform is a classical analytic symbol. The transformed operator $Q=\lambda P$ is then a formal analytic classical pseudodifferential operator of order one, with symbol defined in a complex neighborhood of $(0,0)$.

Consider the vector field
\[
\nu=\sum_{|\alpha|=|\beta|=1}x^\alpha a_{\alpha\beta}(x,0)\partial_x^\beta.
\]
As in Remark~3.10, if $\Omega$ is a suitable small neighborhood of $0$, then $x_t=\exp(t\nu)(x)\in\Omega$ for $t<0$ and
\[
|x_t|\le Ce^{-|t|/C}|x|.
\]
If $f$ is holomorphic in $\Omega$ and $f(0)\notin\Sp(Q_0)$, then for every holomorphic $v$ there is a unique holomorphic $u$ satisfying
\begin{equation}
(\nu-f)u=v\quad\text{in }\Omega.
\end{equation}
Moreover, if $v$ vanishes to order $k$ at $0$, so does $u$. The assertion is immediate for formal power series; existence for holomorphic functions follows by reducing to $v(x)=\cO(|x|^N)$ with $N$ large and using the integral formula (3.27).

Let $\mathcal L^{m,k}(\Omega)$, respectively $\mathcal L_{1/2}^{m,k}(\Omega)$, denote the formal symbols
\begin{equation}
u(x,\lambda)=\sum_{j\in\N\ (j\in\frac12\N)}u_j(x)\lambda^{m-j},
\end{equation}
where $u_j$ is holomorphic in $\Omega$ and vanishes to order $(k-2j)_+$ at the origin. The transport equation just discussed shows that $Q-E$ is bijective on the corresponding formal symbol spaces whenever $E\notin\Sp(Q_0)$.

Let $\mathcal L_a^{m,k}$ and $\mathcal L_{1/2,a}^{m,k}$ denote the subspaces which are formal analytic symbols, i.e. for every compact $K\Subset\Omega$ there is $C>0$ such that
\[
|u_j(x)|\le C^{j+1}j!,\qquad x\in K.
\]
The main step is the following.

\begin{theorem}
For $E\notin\Sp(Q_0)$, the operator $Q-E$ is a bijection
\[
\mathcal L_a^{m,k}(\Omega)\longrightarrow\mathcal L_a^{m,k}(\Omega),
\]
and similarly in the half-order case.
\end{theorem}
The result is related to work of M\'etivier [18]. We give an argument based on weighted holomorphic spaces.

Let $\phi$ be strictly plurisubharmonic and real analytic near $0\in\C^n$, with $\phi(0)=0$, $\phi'(0)=0$. Let $\psi(x,y)$ be the unique holomorphic function near $(0,0)$ with $\psi(x,\bar x)=\phi(x)$. Instead of the standard phase we use
\[
\frac2i\bigl(\psi(x,\theta)-\psi(y,\theta)\bigr).
\]
If $a(x,\theta,\lambda)$ is an analytic symbol near the origin in $\C^{2n}$, consider formally
\begin{equation}
Au(x,\lambda)=\left(\frac{\lambda}{2\pi i}\right)^n
\iint e^{2\lambda(\psi(x,\theta)-\psi(y,\theta))}
 a(x,\theta,\lambda)u(y,\lambda)\,dy\,d\theta.
\end{equation}
Let $\Omega$ be a small strictly pseudoconvex neighborhood of $0$ and let $H_\phi(\Omega)$ be the space of holomorphic functions for which
\[
\norm{u(\cdot,\lambda)}_{H_\phi}\le C_\eps e^{\eps\lambda},\qquad\lambda>0,
\]
where the $L^2$ density is $e^{-2\lambda\phi}L(dx)$. On the contour $\theta=\bar y$ the operator becomes
\begin{equation}
Au(x,\lambda)=\left(\frac\lambda\pi\right)^n
\int_\Omega e^{2\lambda\psi(x,\bar y)}a(x,\bar y,\lambda)
 u(y,\lambda)e^{-2\lambda\phi(y)}L(dy).
\end{equation}
The adjoint is of the same form with $a(x,\bar y,\lambda)$ replaced by the corresponding transposed-conjugate amplitude.

We shall also need complex $\lambda=re^{i\omega}$. Specialize to $\phi(x)=\frac12|x|^2$, so $\psi(x,\bar y)=\frac12x\cdot\bar y$. Choosing the contour $\theta=e^{-i\omega}\bar y$ gives
\begin{equation}
Au(x,\lambda)=\left(\frac r\pi\right)^n
\int e^{r x\cdot\bar y}a(x,e^{-i\omega}\bar y,\lambda)
 u(y,\lambda)e^{-r|y|^2}L(dy),
\end{equation}
and
\begin{equation}
A^*u(x,\lambda)=\left(\frac r\pi\right)^n
\int e^{r y\cdot\bar x}\overline{a(y,e^{-i\omega}\bar x,\lambda)}
 u(y,\lambda)e^{-r|y|^2}L(dy).
\end{equation}
With $r$ as the large parameter, these are ordinary analytic pseudodifferential operators. In the case $A=Q$ the symbol has the form
\begin{equation}
A(x,\xi,\lambda)=i\lambda\sum_{|\alpha|=|\beta|=1}
x^\alpha a_{\alpha\beta}(x,\xi)\xi^\beta+a(x,\xi,\lambda),
\end{equation}
with $a$ of order zero. The symbols of $A$ and $A^*$ are therefore, respectively,
\begin{equation}
ir\sum x^\alpha a_{\alpha\beta}(x,e^{-i\omega}\xi)\xi^\beta
+a(x,e^{-i\omega}\xi,e^{i\omega}r),
\end{equation}
and
\begin{equation}
ir\sum \xi^\alpha a_{\alpha\beta}(\xi,e^{-i\omega}x)x^\beta
+a(\xi,ie^{-i\omega}x,e^{i\omega}r).
\end{equation}
In view of (3.19) and the value of the subprincipal symbol, these agree to the necessary order with the model operator $Q_r$. We therefore write $Q$ and $Q^*$ for these realizations. Let $E\in\C\setminus\Sp(Q_0)$ and put $\mathcal P=Q-E$.

\begin{proposition}
If $\Omega_0\Subset\Omega$ is a neighborhood of $0$, then there is $C>0$ such that
\begin{equation}
\norm v_{H_\phi(\Omega_0)}\le C\left(
\norm{\mathcal P^*v}_{H_\phi(\Omega)}+
\norm v_{H_\phi(\Omega\setminus\Omega_0)}\right),
\end{equation}
for every $v\in H^0_\phi(\Omega)$ and every $|\lambda|\ge1$.
\end{proposition}
We also have
\begin{equation}
\norm{\mathcal P^*v}_{H^0_\phi(\Omega)}
\le C|\lambda|\norm v_{H^0_\phi(\Omega)},
\end{equation}
and the same type of estimate for $\nabla_\lambda\mathcal P^*$.

Consequently the operator
\[
\mathcal A^*:v\longmapsto\bigl(\mathcal P^*v,\ v|_{\Omega\setminus\Omega_0}\bigr)
\]
is injective with closed range. Let $\mathcal B=(\mathcal A^*\mathcal A)^{-1}\mathcal A^*$ be its optimal left inverse. Then
\[
\norm{\mathcal B}=\cO(|\lambda|),
\qquad
\norm{\nabla_\lambda\mathcal B}=\cO(|\lambda|^3).
\]
To justify differentiation of adjoints, note that the weighted scalar products depend on $r=|\lambda|$. If $A_\lambda$ is a $C^1$ family and $A^\#$ its adjoint for the $r$-dependent scalar products, then
\[
\partial_rA^\#=(FA-AF+\partial_rA)^\#,
\qquad
\partial_\omega A^\#=(\partial_\omega A)^\#,
\]
with $\norm F=\cO(1)$; hence
\[
\norm{\nabla_\lambda A^\#}=\cO(1)
\bigl(\norm A+\norm{\nabla_\lambda A}\bigr).
\]
The weighted Bergman projection $B_{\phi,\Omega}$ satisfies $\norm{\nabla_\lambda B_{\phi,\Omega}}=\cO(1)$.

Taking adjoints of the left-inverse relation gives operators $C,D$ such that
\begin{equation}
I=\mathcal P C+B_{\phi,\Omega}\,1_{\Omega\setminus\Omega_0}D.
\end{equation}
Thus
\[
\mathcal P Cv=v-Gv,
\qquad Gv:=B_{\phi,\Omega}(1_{\Omega\setminus\Omega_0}Dv).
\]
We next estimate the Bergman projection. The identity on $H_\phi(\Omega)$ has a realization
\begin{equation}
Jg(x,\lambda)=\left(\frac r\pi\right)^n
\int e^{2r\psi(x,\bar y)}a(x,\bar y,r)g(y,\lambda)e^{-2r\phi(y)}L(dy),
\end{equation}
where $a$ is a classical analytic symbol of order zero. If $C_0$ is large and
\[
\widetilde\phi(x)=\phi(x)-\frac{d(x,\complement\Omega)^2}{C_0},
\]
then
\begin{equation}
\norm{g-Jg}_{H_{\widetilde\phi}(\Omega)}
\le C\norm g_{H_\phi(\Omega)},\qquad r\ge1.
\end{equation}
Since $Jg=Jf$ when $g=B_{\phi,\Omega}f$,
\begin{equation}
\norm{B_{\phi,\Omega}f-Jf}_{H_{\widetilde\phi}(\Omega)}
\le C\norm f_{L^2_\phi(\Omega)}.
\end{equation}
It follows that, after slightly decreasing the weight on $\Omega_0$,
\[
G=\cO(|\lambda|),
\qquad
\nabla_\lambda G=\cO(|\lambda|^3):
H_\phi(\Omega)\to H_{\widetilde\phi}(\Omega_0).
\]

Let $v_f=\sum_{k\ge0}v_k(x)\lambda^{-k}$ be a formal analytic symbol. Choose a realization
\[
v(x,\lambda)=\sum_kv_k(x)\lambda^{-k}\chi(k/(C\lambda)),\qquad|\lambda|\ge1,
\]
with $C$ large. Then $v=\cO(1)$ and
\begin{equation}
\partial_\lambda v=\cO(e^{-C|\lambda|})
\end{equation}
in the sup norm on $\Omega$. Put $u=Cv$, so that
\begin{equation}
\mathcal Pu=v-Gv\quad\text{in }\Omega_0,
\qquad u=\cO(|\lambda|)\text{ in }H^0_\phi(\Omega).
\end{equation}
Differentiating gives
\begin{equation}
\mathcal P(\partial_\lambda u)
=\partial_\lambda v-(\partial_\lambda G)v-G\partial_\lambda v
-(\partial_\lambda\mathcal P)u=:w.
\end{equation}
The first three terms satisfy
\begin{equation}
(\partial_\lambda G)v-G\partial_\lambda v
=\cO(|\lambda|^3)\quad\text{in }H_{\widetilde\phi}(\Omega_0).
\end{equation}
Using Stokes' formula on the contour realization shows that $\partial_\lambda\mathcal P$ is $\cO(1)$ from $H_\phi(\Omega)$ to a slightly weaker weighted space, and hence
\begin{equation}
w=\cO(|\lambda|^3)\quad\text{in }H_{\widetilde\phi}(\Omega_0).
\end{equation}
The a priori estimate is stable under small perturbations $\phi_\eps=\phi+\eps|x|^2$. If $\Omega_1\Subset\Omega_0$,
\begin{equation}
\norm{\partial_\lambda u}_{H_{\phi_\eps}(\Omega_1)}
\le C\left(
\norm w_{H_{\widetilde\phi}(\Omega_0)}+
\norm{\partial_\lambda u}_{H_{\phi_\eps}(\Omega_0\setminus\Omega_1)}
\right).
\end{equation}
The right hand side is exponentially decreasing after choosing the weights suitably. In particular,
\begin{equation}
\norm{\partial_\lambda u}_{H_{\widetilde\phi}(\Omega_0)}=\cO(1),
\qquad |\lambda|\ge1.
\end{equation}
We extend $u$ smoothly to $|\lambda|<1$ while keeping it holomorphic in $x$.

\begin{proposition}
Let $v_f$ be the formal analytic symbol corresponding to $v$, and let
\[
u_f=\sum_j u_j(x)\lambda^{-j}
\]
be the unique formal symbol satisfying $\mathcal P u_f=v_f$. Then $u_f$ is a formal analytic symbol, and $u$ is a realization of $u_f$ near $0$.
\end{proposition}
\begin{proof}
Let $u_f^N=\sum_{j=0}^N u_j(x)\lambda^{-j}$. Then
\[
\mathcal P u_f^N=v+\cO_\Omega(\lambda^{-N})
\]
uniformly in $\Omega_0$. The a priori estimate with $\mathcal P$ in place of $\mathcal P^*$ gives $u_f^N-u=\cO(\lambda^{-N})$ in a smaller weighted space. Hence for every multi-index $\alpha$,
\begin{equation}
\partial_x^\alpha u(0,\lambda)=
\sum_{j=0}^N\partial_x^\alpha u_j(0)\lambda^{-j}
+\cO_{\alpha,N}(\lambda^{-N-1}).
\end{equation}
The following lemma then yields the analytic growth of the coefficients.
\end{proof}

\begin{lemma}
After decreasing $\Omega_1$ if necessary,
\begin{equation}
u_N(x)=\lim_{R\to\infty}\frac1{2\pi i}
\int_{|\lambda|=R}\lambda^{N-1}u(x,\lambda)\,d\lambda
=\frac1{2\pi i}\iint_\Gamma
\lambda^{N-1}\frac{\partial u}{\partial\bar\lambda}
\,d\bar\lambda\wedge d\lambda.
\end{equation}
\end{lemma}
\begin{proof}
The exponential decrease of $\partial_{\bar\lambda}u$ makes the last integral converge and Stokes' formula gives the second identity. The first identity follows from the asymptotic expansion at $x=0$ and then, by holomorphy, throughout a smaller neighborhood.
\end{proof}
The lemma implies
\[
|u_N(x)|\le C\iint |\lambda|^{N-1}e^{-C_0|\lambda|}
|d\lambda\wedge d\bar\lambda|
\le C_1C_0^{-N}N!,
\]
so $u_f$ is a formal analytic symbol. Any standard analytic realization differs from $u$ by an exponentially small term. This completes the difficult part of Theorem~4.1; away from the origin, analyticity propagates along the integral curves of $\nu$ by the principal-type theory of [24]. The half-order case is handled by separating the two parity series.

We now return to the spectral projection $\Pi$ of Section~3. The basis elements $e_j$ are analytic symbols, and the uniformity in Theorem~4.1 implies that $\Pi$ preserves analytic symbols. Stationary phase preserves analytic symbols, so the Gram matrix and its inverse square root are analytic symbols; hence the orthonormal basis $f_1,\ldots,f_{N_0}$ and the finite Hermitian matrix $Q$ in (3.26) are analytic symbols.

Let $\lambda_1\in S_{1/2}^{-1/2}(\R_+)$ be a formal eigenvalue of multiplicity $d$. Choose a smooth double-valued realization $\widetilde\lambda_1$ and a realization $\widetilde Q$ with $\partial_\lambda\widetilde Q=\cO(e^{-|\lambda|/C})$, Hermitian for $\lambda>0$. For $N_1$ sufficiently large, the spectral projection onto the $d$ eigenvalues in
$|E-\widetilde\lambda_1|\le|\lambda|^{-N_1}$ is
\[
\widetilde\Pi_1=\frac1{2\pi i}
\int_{|E-\widetilde\lambda_1|=|\lambda|^{-N_1}}
(E-\widetilde Q)^{-1}\,dE.
\]
It has an asymptotic expansion $\Pi_1$ of order zero and $\partial_\lambda\widetilde\Pi_1=\cO(e^{-|\lambda|/C})$. By the same argument as above, $\Pi_1$ is an analytic symbol. Since
\[
\lambda_1=d^{-1}\operatorname{tr}(\Pi_1Q),
\]
the formal eigenvalue is an analytic symbol, and similarly for the other eigenvalues and for an orthonormal eigenbasis.

\begin{theorem}
The $g_j$ and the corresponding eigenvalues in Theorem~3.6 can be chosen to be analytic symbols.
\end{theorem}

The calculations above used the weight $\frac12|z|^2$ rather than $\frac14|z|^2$, but an elementary change of variables connects the two. Since FBI transforms and their inverses preserve analytic symbols, we obtain:

\begin{theorem}
The amplitudes $a_j$ and eigenvalues $E_j$ in Theorem~3.7 can be chosen to be analytic symbols. If $\widetilde a_j$ and $\widetilde E_j$ are realizations and $\widetilde u_j=\widetilde a_je^{-\lambda\psi}$, then there is $C>0$ such that
\begin{equation}
\lambda P\widetilde u_j-\widetilde E_j\widetilde u_j
=\cO(e^{-\lambda/C}),
\end{equation}
uniformly for $\lambda\ge1$ and $|x|\le1/C$.
\end{theorem}
The amplitudes and realizations may be taken real valued. In one dimension F. Pham independently proved this result; the source cites his Strasbourg R.C.P. seminar contribution of May 1983.


\section{The interaction matrix in the case of strict minima}
The first part of this section is devoted to various estimates for a fixed $M_j=M_1$ under the same assumptions as in Section~3, sometimes also with the analyticity assumptions. The first task is to show that the asymptotic eigenvalues of the preceding sections are close to the true eigenvalues and similarly for the eigenfunctions in $L^2$ norm. After that we show, under a suitable additional hypothesis, that the approximate eigenfunctions are close to the true ones far away from $U_j=\{0\}$, even when the exponential decrease factor is taken into account.

Let $R_0>0$ be fixed and avoid $\Sp(Q_0)$. We consider eigenvalues $E$ in
\[
I(h)=[0,R_0h]
\]
(or occasionally in a larger interval of the same type). If $\Omega\Subset M_1$ is a sufficiently small coordinate neighborhood of $0$, we have the following estimate.

\begin{proposition}
For every $M\in\N$ there is a constant $C_M>0$ such that, for all $E\in I(h)$, $h>0$ sufficiently small and $u\in C_0^\infty(\Omega)$,
\begin{equation}
\norm{u}_{M+2}\le C_M\left(
\norm{\frac1h(P-E)u}_M+\norm{u}_0\right).
\tag{5.1}
\end{equation}
Here
\begin{equation}
\norm{u}_M=
\sum_{|\alpha|+|\beta|\le M}
 h^{-(|\alpha|+|\beta|)/2}
 \norm{x^\alpha(hD)^\beta u}_{L^2(\Omega)}.
\tag{5.2}
\end{equation}
\end{proposition}
\begin{proof}
Write
\[
P=-h^2\Delta+V
 =\sum_{|\alpha|=2}a_\alpha(x)(hD)^\alpha
 +h\sum_{|\alpha|=1}a_\alpha(x)(hD)^\alpha+V(x)
\]
and put
\[
P_0=\sum_{|\alpha|=2}a_\alpha(0)(hD)^\alpha
+\frac12\langle V''(0)x,x\rangle.
\]
Then
\[
P-P_0=\cO(x)(hD)^2+\cO(1)h(hD)+\cO(x^3).
\]
The operator $P_0$ is a harmonic oscillator, unitarily equivalent to $hQ_0$, and (5.1) is standard for $P_0$ when $M=0$. For each $\eps>0$, after decreasing $\Omega$,
\[
\frac1h\norm{(P-P_0)u}\le\eps\norm{u}_2,
\]
so (5.1) also holds for $P$ when $M=0$.

Assume inductively that (5.1) is known for all indices up to $M$. We use
\[
\norm{u}_{M+3}\le C\norm{u}_{M+2}
 +\sum_{|\alpha|=1}h^{-1/2}\norm{x^\alpha u}_{M+2}
 +\sum_{|\beta|=1}h^{-1/2}\norm{hD^\beta u}_{M+2}.
\]
For the first sum the induction hypothesis gives
\begin{equation}
 h^{-1/2}\norm{x^\alpha u}_{M+2}
 \le C\left(
 h^{-1/2}\norm{\frac1h(P-E)x^\alpha u}_M
 +h^{-1/2}\norm{x^\alpha u}_0\right).
\tag{5.3}
\end{equation}
The last term is controlled by $\norm{u}_2$, hence by the right hand side of (5.1) for $M=0$. The first term on the right of (5.3) is bounded by
\begin{equation}
 h^{-1/2}\norm{x^\alpha\frac1h(P-E)u}_M
 +h^{-3/2}\norm{[P,x^\alpha]u}_M.
\tag{5.4}
\end{equation}
Since
\[
[P,x^\alpha]=\cO(1)h(hD)+\cO(1)h^2,
\]
the second term in (5.4) is bounded by
\[
C\left(h^{-1/2}\norm{hDu}_M+h^{1/2}\norm{u}_M\right)
\le C'\norm{u}_{M+2}.
\]
The terms with $hD^\beta u$ are treated in the same way, completing the induction.
\end{proof}

Let $u$ be a normalized eigenfunction of $P_{M_1}$ with eigenvalue $E\in I(h)$. Choose $\chi\in C_0^\infty(\Omega)$ equal to $1$ near $0$ and put $v=\chi u$. The decay results of Section~2, together with ordinary elliptic estimates, show that all derivatives of $u$ are exponentially small on compact subsets of $\Omega\setminus\{0\}$. Consequently $\norm{u-v}_{L^2}$ is exponentially small, $\norm v\sim1$, and
\[
Pv=Ev+[P,\chi]u,
\qquad
\norm{D^\alpha[P,\chi]u}\le C_\alpha e^{-C/h}.
\]
Proposition~5.1 gives $\norm v_M\le C_M$ for every $M$. Since
\[
\norm{(P-P_0)v}\le C h^{3/2}\norm v_3,
\]
we obtain
\begin{equation}
P_0v=Ev+w,
\qquad w=\cO(h^{3/2})\quad\hbox{in }L^2.
\tag{5.5}
\end{equation}
Thus, for some $C_0>0$ and $h$ small,
\begin{equation}
\Sp(P_{M_1})\cap I(h)
\subset \Sp(P_0)+B(0,C_0h^{3/2}).
\tag{5.6}
\end{equation}
Recall that $\Sp(P_0)=h\Sp(Q_0)$ is explicit. By deforming $P$ and the metric inside $\Omega$ to the harmonic oscillator model, without changing anything outside $\Omega$, the arguments of Section~2 show that the multiplicities are preserved. Hence:

\begin{proposition}
There is a bijection
\[
b:\Sp(P_{M_1})\cap I(h)\longrightarrow \Sp(P_0)\cap I(h)
\]
such that $|b(E)-E|\le C_0h^{3/2}$ for $h$ sufficiently small.
\end{proposition}

Let $E_0\in\Sp(Q_0)$ have multiplicity $N_0$ and let $I_{E_0}(h)$ be the corresponding cluster of eigenvalues.

\begin{proposition}
Let $\widehat a_j,\widehat E_j$ be the formal quantities of Remark~3.10 and put
\[
\widehat u_j=h^{-n/4}\widehat a_j e^{-\phi/h}\chi,
\qquad j=1,\ldots,N_0.
\]
If $\widehat E_1$ denotes the space spanned by the $\widehat u_j$, then
\[
\dvec(\widehat E_1,E_1)=\cO(h^\infty),
\]
and the eigenvalues of $P_{M_1}$ in $I_{E_0}(h)$ are
$\widehat E_j+\cO(h^\infty)$.
\end{proposition}

\begin{proposition}
In the analytic case, let $\widehat a_j,\widehat E_j$ be realizations as in Theorem~4.6 and let $\widehat E_1$ be spanned by
\[
\widehat u_j=h^{-n/4}\widehat a_j e^{-\phi/h}\chi.
\]
Then, for some $C>0$,
\[
\dvec(\widehat E_1,E_1)=\cO(e^{-1/(Ch)}),
\]
and the eigenvalues of $P_{M_1}$ in $I_{E_0}(h)$ are
$\widehat E_j+\cO(e^{-1/(Ch)})$.
\end{proposition}

We next study eigenfunctions far away from the well. Let $d(x,y)$ denote the Agmon distance and write $d(x)=d(x,0)$.

\begin{proposition}
There exist constants $C,\widetilde C>0$ such that every normalized eigenfunction $u$ of $P_{M_1}$ with eigenvalue $E\in I(h)$ satisfies
\begin{equation}
 h^2\norm{(1+d/h)^{-C}e^{d/h}\nabla u}^2
 +h\norm{(1+d/h)^{-C}e^{d/h}u}^2
 \le \widetilde C h.
\tag{5.7}
\end{equation}
\end{proposition}
\begin{proof}
Choose
\[
\phi(x)=
\begin{cases}
 d(x)-Ch\log(d(x)/h),&d(x)\ge Ch,\\
 d(x)-Ch\log C,&d(x)\le Ch.
\end{cases}
\]
For $d(x)\le Ch$, the distance is smooth and $V-|\nabla\phi|^2=0$. For $d(x)\ge Ch$,
\[
\nabla\phi=\left(1-\frac{Ch}{d(x)}\right)\nabla d,
\]
so
\[
V(x)-|\nabla\phi(x)|^2\ge C h\frac{V(x)}{d(x)}.
\]
Near the non-degenerate minimum there is $C_0>0$ such that
$C_0^{-1}\le V(x)/d(x)\le C_0$, and therefore
\[
V-|\nabla\phi|^2\ge \frac{C}{C_0}h
\qquad(d\ge Ch).
\]
Apply (1.3) of Theorem~1.1 to $P-E$. We may take
\[
F_+=\mathbf1_{\{d\ge Ch\}}
 \sqrt{V-E-|\nabla\phi|^2}
\ge \mathbf1_{\{d\ge Ch\}}
 \left(\frac{C}{C_0}h-E\right)^{1/2},
\quad
F_-=\mathbf1_{\{d\le Ch\}}E^{1/2},
\]
choosing $C/C_0>E/h+1$. This yields
\begin{equation}
 h^2\norm{\nabla(e^{\phi/h}u)}^2
 +\frac12\left(\frac{C}{C_0}h-E\right)
 \norm{e^{\phi/h}u}_{\{d\ge Ch\}}^2
 \le \frac32 E\norm{e^{\phi/h}u}_{\{d\le Ch\}}^2.
\tag{5.8}
\end{equation}
The weight $e^{\phi/h}$ is comparable with $(1+d/h)^{-C}e^{d/h}$, and the derivative of the weight is controlled by $h^{-1/2}(d/h)^{1/2}e^{\phi/h}$. Equation (5.7) follows.
\end{proof}

Let $u_1,\ldots,u_{N_0}$ be an orthonormal basis of exact eigenfunctions and let $E_1,\ldots,E_{N_0}$ be the corresponding eigenvalues. Propositions~5.3 and~5.4 imply that an orthogonal matrix $(c_{j,k}(h))$ can be chosen so that
\begin{equation}
u_j=\sum_k c_{j,k}\widehat u_k+\cO(h^\infty)
\quad\hbox{in }L^2(M_1),
\tag{5.9}
\end{equation}
and, in the analytic case,
\begin{equation}
u_j=\sum_k c_{j,k}\widehat u_k+\cO(e^{-1/(Ch)}).
\tag{5.10}
\end{equation}
Moreover $c_{j,k}$ may be taken equal to zero when the asymptotic eigenvalue attached to $\widehat u_k$ is not asymptotically equal to $E_j$. In particular, if all the formal eigenvalues have different asymptotic expansions, the matrix $(c_{j,k})$ may be chosen to be the identity.

Let $\Omega\subset M_1$ be open, contain $0$, and assume
\begin{equation}
d(x)\quad\hbox{is a }C^2\hbox{ function on }\Omega.
\tag{5.11}
\end{equation}
Then
\[
\Lambda_+=\{(x,\nabla d(x));\ x\in\Omega\}\subset q^{-1}(0),
\qquad q(x,\xi)=\xi^2-V(x),
\]
and $H_q$ is tangent to $\Lambda_+$. We also assume
\begin{equation}
(x,\xi)\in\Lambda_+\Longrightarrow
\exp(tH_q)(x,\xi)\in\Lambda_+\ (t\le0),
\quad
\exp(tH_q)(x,\xi)\to(0,0)\ (t\to-\infty).
\tag{5.12}
\end{equation}
Since $\Lambda_+$ is $C^\infty$ (respectively analytic) near $(0,0)$, it has the same regularity everywhere along these trajectories, and hence
\begin{equation}
d(x)\quad\hbox{is }C^\infty\quad\hbox{(respectively analytic) on }\Omega.
\tag{5.13}
\end{equation}
As in Section~3, $d(x)$ is the length, in the Agmon metric, of the projected trajectory from $0$ to $x$.

From now on all wells are assumed to be points $U_j=\{x_j\}$, and
\begin{equation}
V''(x_j)>0,
\qquad 1\le j\le N.
\tag{5.14}
\end{equation}

\begin{lemma}
Let $x,y\in\Omega$ and suppose that $y$ does not lie on the minimal segment from $0$ to $x$ obtained by projecting the backward $H_q$-trajectory. Then
\[
d(0,x)<d(0,y)+d(y,x).
\]
\end{lemma}
\begin{proof}
Equip $M$ with the Agmon metric $V(x)dx^2$. The projected Hamilton trajectory from $x$ to $0$ is a minimizing geodesic. It is in fact the unique minimizing geodesic, up to reparametrization: since $d$ is $C^1$ near $x$ and $d(x)\ne0$, every minimizing curve issuing from $x$ has tangent proportional to the negative gradient of $d$ and hence coincides with this trajectory. The triangle inequality gives $d(0,x)\le d(0,y)+d(y,x)$. Equality would concatenate minimizing segments into a second minimizing geodesic from $x$ to $0$, contrary to uniqueness.
\end{proof}

\begin{lemma}
Let $K_1,K_2\subset\Omega$ be compact and suppose that $K_2$ is disjoint from the compact union $\widehat K_1$ of all minimal geodesics from points of $K_1$ to $0$. Then there is $\eps>0$ such that, for all $x\in K_1$, $y\in K_2$,
\begin{equation}
d(0,x)\le(1-\eps)\bigl(d(0,y)+d(y,x)\bigr).
\tag{5.15}
\end{equation}
\end{lemma}

\begin{theorem}
Let $\Omega\Subset M_1$ satisfy (5.11) and (5.12). Put $\psi=d(x)$. Then the various amplitudes of Remark~3.10, Theorem~4.6 and Propositions~5.3--5.4 may be defined throughout $\Omega$ by solving the transport equations globally. Thus the approximate functions $\widehat u_j$ in (5.9), (5.10) are defined there as well. If
\[
u'_j=\sum_k c_{j,k}\widehat u_k,
\]
then:

\begin{enumerate}[label=(\Alph*)]
\item For every compact $K\Subset\Omega$ and every $N\in\N$,
\begin{equation}
\norm{\nabla\bigl(e^{d/h}(u_j-u'_j)\bigr)}_{L^2(K)}^2
+\norm{e^{d/h}(u_j-u'_j)}_{L^2(K)}^2
=\cO(h^N).
\tag{5.16}
\end{equation}
\item In the analytic case the right hand side of (5.16) may be replaced by $\cO(e^{-1/(Ch)})$, where $C>0$ may depend on $K$.
\end{enumerate}
\end{theorem}
\begin{proof}
Set $w=u_j-u'_j$. Proposition~5.5 and the explicit form of $u'_j$ show first that the left hand side of (5.16) grows at most polynomially in $h^{-1}$. Moreover
\begin{equation}
Pw=E_jw+r,
\tag{5.17}
\end{equation}
where $r=(P-E_j)u'_j$ satisfies, for every compact $K\Subset\Omega$ and every $N$,
\begin{equation}
\norm{e^{d/h}r}_{L^2(K)}=\cO(h^N),
\tag{5.18}
\end{equation}
and (5.9) gives
\begin{equation}
\norm{w}_{L^2(K)}=\cO(h^N).
\tag{5.19}
\end{equation}
Fix $K$ and choose $\chi\in C_0^\infty(\Omega)$ equal to $1$ in a neighborhood of the union $\widehat K$ of all minimal geodesics from $K$ to $0$. If $\phi$ is the weight used in Proposition~5.5, apply Theorem~1.1 to $\chi w$ with the new weight
\begin{equation}
\phi_N(x)=\min\left(
\phi(x)+Nh\log(1/h),
\inf_{y\in\supp\nabla\chi}\bigl(\phi(y)+(1-\eps)d(y,x)\bigr)
\right).
\tag{5.20}
\end{equation}
By Lemma~5.7, for $\eps$ small and $h<h_N$, $\phi_N=\phi+Nh\log(1/h)$ near $\widehat K$. On the other hand, where the second branch is active,
\[
V-|\nabla\phi_N|^2\ge(2\eps-\eps^2)V.
\]
Combining this with the estimate used in Proposition~5.5 gives, almost everywhere,
\begin{equation}
V(x)-|\nabla\phi_N(x)|^2\ge
\begin{cases}
\eps Ch/C_0,& d(x)\ge Ch,\\
0,&d(x)\le Ch.
\end{cases}
\tag{5.21}
\end{equation}
With $C/C_0>E_j/h+1$, Theorem~1.1 yields an estimate of the form
\begin{align}
&h^{2-2N}\norm{\nabla(e^{\phi_N/h}w)}_{L^2(K)}^2
+h^{1-2N}\norm{e^{\phi_N/h}w}_{L^2(K)}^2 \notag\\
&\qquad\le C_1\left(
 h^{1-2N}\norm{e^{\phi_N/h}w}_{L^2(\{d\le Ch\})}^2
 +h^{-1}\norm{e^{\phi_N/h}(P-E_j)\chi w}^2
\right).
\tag{5.22}
\end{align}
The first term on the right is bounded by (5.19). For the last term,
\begin{align}
&2h^{-1}\norm{e^{\phi_N/h}\chi(P-E_j)w}^2
+2h^{-1}\norm{e^{\phi_N/h}[P,\chi]w}^2 
otag\\
&\qquad\le
2h^{-1}(1+2N)\norm{e^{\phi/h}\chi(P-E_j)w}^2
+2h^{-1}\norm{e^{\phi/h}[P,\chi]w}^2,
\tag{5.23}
\end{align}
where we used $\phi_N\le\phi$ on $\supp\nabla\chi$. Equations (5.17), (5.18) bound the first term and the second has at most a fixed polynomial growth. Since $N$ is arbitrary, (A) follows.

In the analytic case, (5.18), (5.19) have exponentially small right hand sides. Replace $\phi_N$ by
\begin{equation}
\widetilde\phi(x)=(1-\eps)\min\left(
 d(x)+a,
 \inf_{y\in\supp\nabla\chi}\bigl(d(y)+d(y,x)\bigr)
\right).
\tag{5.24}
\end{equation}
For $a>0$ small, $\widetilde\phi=(1-\eps)(d+a)$ near $\widehat K$, while $V-|\nabla\widetilde\phi|^2\ge c\eps$ away from the well. The weighted identity gives
\begin{align}
&\left(
 h^2\norm{\nabla(e^{(1-\eps)d/h}w)}_{L^2(K)}^2
 +\norm{e^{(1-\eps)d/h}w}_{L^2(K)}^2
 \right)e^{2(1-\eps)a/h} \notag\\
&\quad\le C_{r,\eps,a}\left(
 \norm{e^{(1-\eps)(d+a)/h}w}_{L^2(\{d<r\})}^2
 +\norm{e^{\widetilde\phi/h}\chi(P-E_j)w}^2
 +\norm{e^{\widetilde\phi/h}[P,\chi]w}^2
\right).
\tag{5.25}
\end{align}
Choosing $a,r$ sufficiently small makes the right hand side bounded uniformly as $h\to0$; then choosing $\eps>0$ sufficiently small proves (B).
\end{proof}

As before, let
\[
S_0=\min_{j\ne k}d(x_j,x_k).
\]
Choose open neighborhoods $\Omega_j\Subset M$ of $x_j$, with $x_k\notin\overline\Omega_j$ for $k\ne j$, satisfying (5.11), (5.12), with $d_j(x)=d(x,x_j)$. A geodesic for the Agmon metric joins any two points and is $C^\infty$ away from the wells. We now assume:
\begin{equation}
 d(x_j,x_k)=S_0\Longrightarrow
 \text{there is a unique minimizing geodesic }\gamma_{j,k}
 \text{ joining }x_j\text{ to }x_k.
\tag{5.26}
\end{equation}
\begin{equation}
\gamma_{j,k}\subset\Omega_j\cup\Omega_k.
\tag{5.27}
\end{equation}
If $x_0\in\gamma_{j,k}$, $\Omega_j\cap\Omega_k\ne\varnothing$, and $\Gamma\subset\Omega_j\cap\Omega_k$ is a smooth hypersurface meeting $\gamma_{j,k}$ transversally at $x_0$ and nowhere else, assume that there is $C>0$ such that
\begin{equation}
d(x,x_j)+d(x,x_k)
\ge S_0+\frac1C d(x,x_0)^2,
\qquad x\in\Gamma.
\tag{5.28}
\end{equation}
This condition is independent of the particular transverse section. We choose $M_j$ with $S>S_0/2$ and $\Omega_j\Subset\mathring M_j$.

Let $E_0$ be in the spectrum of some localized harmonic oscillator and let
$E(h)=E_0h+\cO(h^{3/2})$ be one of the asymptotic eigenvalues constructed above. Choose
\[
I(h)=[E(h)-h^{N_0},E(h)+h^{N_0}]
\]
with $N_0$ sufficiently large. Then the localized eigenvalues $\mu_\alpha$ in $I(h)$ agree with $E(h)$ to all orders. We shall assume, for convenience,
\begin{equation}
j(\alpha)=j(\beta)\Longrightarrow\mu_\alpha=\mu_\beta.
\tag{5.29}
\end{equation}
In the analytic case there is $b>0$ such that
\begin{equation}
|\mu_\alpha-\mu_\beta|=\cO(e^{-b/h}),
\qquad\forall\alpha,\beta.
\tag{5.30}
\end{equation}
Let $\Gamma$ be as in (5.28). By Remark~2.13, when $d(x_{j(\alpha)},x_{j(\beta)})=S_0$,
\begin{equation}
\widehat w_{\alpha,\beta}
=h^2\int_\Gamma
\left(
\varphi_\alpha\frac{\partial\varphi_\beta}{\partial n}
-\varphi_\beta\frac{\partial\varphi_\alpha}{\partial n}
\right)dS,
\tag{5.31}
\end{equation}
where $n$ is the unit normal oriented towards $x_{j(\alpha)}$. Equality here is modulo $\cO(h^N e^{-S_0/h})$ for every $N$ in the $C^\infty$ case and modulo $\cO(e^{-(S_0+a)/h})$ for some $a>0$ in the analytic case.

Applying Theorem~5.8 and substituting the WKB expressions into (5.31), we get
\begin{equation}
\widehat w_{\alpha,\beta}
=h^{1-n/2}\int_\Gamma
 a_{\alpha,\beta}
 e^{-(\psi_{j(\alpha)}+\psi_{j(\beta)})/h}\,dS,
\tag{5.32}
\end{equation}
where
\begin{equation}
a_{\alpha,\beta}
=a_\alpha a_\beta\left(
\frac{\partial\psi_\alpha}{\partial n}
-\frac{\partial\psi_\beta}{\partial n}
\right)
+h\left(
 a_\alpha\frac{\partial a_\beta}{\partial n}
-a_\beta\frac{\partial a_\alpha}{\partial n}
\right).
\tag{5.33}
\end{equation}
By (5.28), stationary phase applies and gives
\begin{equation}
\widehat w_{\alpha,\beta}=b_{\alpha,\beta}(h)e^{-S_0/h},
\qquad b_{\alpha,\beta}=b_{\beta,\alpha}.
\tag{5.34}
\end{equation}
The coefficient is a classical symbol (an analytic symbol in the analytic case), with computable order.

\begin{theorem}
Choose $I(h)$ as above and assume (5.14), (5.26)--(5.29). Let $\widehat w_{\alpha,\beta}$ be the coefficients of Theorem~2.10. If $d(x_{j(\alpha)},x_{j(\beta)})=S_0$, then they are given by (5.34) up to an error which, in the $C^\infty$ case, is $\cO(h^N e^{-S_0/h})$ for every $N\in\N$ and, in the analytic case, is $\cO(e^{-(S_0+a)/h})$ for some $a>0$.

In the two cases, $b_{\alpha,\beta}$ is respectively a classical or classical analytic symbol of the form
\begin{equation}
b_{\alpha,\beta}(\lambda)
\sim \lambda^{m(\alpha,\beta)}
\sum_{j\in\frac12\N}b_{\alpha,\beta}^{(j)}\lambda^{-j},
\qquad \lambda=h^{-1}.
\tag{5.35}
\end{equation}
The order $m(\alpha,\beta)$ can be estimated explicitly. If $E_0$ can occur in the spectra of the localized harmonic oscillators only as the principal eigenvalue, then all the $b_{\alpha,\beta}$ are elliptic symbols, strictly negative, and $m(\alpha,\beta)=-\frac12$; in that case the sum in (5.35) runs over $\N$.
\end{theorem}

The leading symbols can be computed directly from the transport equations; in this way one also checks that the leading term does not depend on the choice of $\Gamma$.

\begin{remark}
When $P$ is invariant under a group of isometries, the interaction matrix has the corresponding invariance property and the eigenvalue splitting can often be studied explicitly.
\end{remark}

\begin{remark}
In the analytic case the remainder estimates are stronger, making it possible to study small perturbations of symmetric situations. In particular, after modifying the hypotheses, one may hope to describe $\widehat w_{\alpha,\beta}$ when the distance between the corresponding wells is larger than, but very close to, $S_0$.
\end{remark}

\begin{remark}
In one dimension it is not necessary to remain near the bottom of the wells to obtain explicit results. It suffices to assume that, on the relevant energy surface, the turning points are simple; then the $\varphi_\alpha$ can be computed by WKB or Maslov methods. We have developed the multidimensional case first because the geometry is more substantial.
\end{remark}

\section{The case when (5.26)--(5.28) are not necessarily satisfied}
In the first part of this section we recall a result on the non-degenerate behaviour near the interior of a minimal geodesic. In a slightly different form this result is well known and was communicated to the authors by Prof.~Kobayashi. We give a proof using only standard Riemannian geometry, as in Milnor [19].

Let $M$ be a complete Riemannian manifold and $N\subset M$ a compact submanifold. Consider a geodesic $\gamma:[0,1]\to M$ with $\gamma(0)\in N$ and $\dot\gamma(0)\in T_N^\perp$. Let
\[
\Gamma:[0,1]\times(-\eps,\eps)^2\longrightarrow M
\]
be a variation of $\gamma$: $\Gamma(t,0,0)=\gamma(t)$, the required mixed derivatives exist and are continuous, and $\Gamma$ is smooth on the strips of a finite subdivision of $[0,1]$. Put
\begin{equation}
E(\Gamma)=E(\Gamma)(u,v)
=\frac12\int_0^1
\left|\frac{\partial\Gamma}{\partial t}(t,u,v)\right|^2dt.
\tag{6.1}
\end{equation}

\begin{proposition}
Let $\gamma$ be as above and let $\Gamma$ be a variation with $\Gamma|_{t=0}\in N$ and $\Gamma|_{t=1}=\gamma(1)$. At $u=v=0$,
\begin{equation}
\frac{\partial E}{\partial u}=0,
\tag{6.2}
\end{equation}
and
\begin{align}
\frac{\partial^2E}{\partial u\partial v}
={}&\int_0^1
\left\langle
\frac{D}{\partial t}\frac{\partial\Gamma}{\partial u},
\frac{D}{\partial t}\frac{\partial\Gamma}{\partial v}
\right\rangle dt \\notag
&+\int_0^1
\left\langle
R\left(\frac{\partial\Gamma}{\partial u},\dot\gamma\right)\dot\gamma,
\frac{\partial\Gamma}{\partial v}
\right\rangle dt
-\left\langle
\dot\gamma(0),
\frac{D}{\partial u}\frac{\partial\Gamma}{\partial v}
\right\rangle_{t=0}.
\tag{6.3}
\end{align}
\end{proposition}
Here $D/\partial t$ denotes covariant differentiation along $\gamma$ and $R$ is the curvature tensor. The right hand side is a symmetric bilinear form in the piecewise smooth vector fields
$U=\partial_u\Gamma$ and $V=\partial_v\Gamma$. The boundary term is a symmetric bilinear function of $U(0)$ and $V(0)$.

Replacing $[0,1]$ by $[0,\theta]$, $0<\theta<1$, write $E_{[0,\theta]}$ for the corresponding energy and
\begin{equation}
Q_\theta(U,V)=
\int_0^\theta\left\langle\frac{D U}{\partial t},\frac{D V}{\partial t}\right\rangle dt
+\int_0^\theta\langle R(U,\dot\gamma)\dot\gamma,V\rangle dt
+A(U(0),V(0)).
\tag{6.4}
\end{equation}
From now on consider only variations satisfying the usual orthogonality condition at $N$ and the fixed endpoint condition. Choose local coordinates $x=(x',x'')$ near $\gamma(0)$ so that $N$ is given by $x''=0$ and $T_N^\perp$ by the $\partial_{x''}$ directions. Then the admissible vector fields satisfy
\[
\pi_{x''}U(0)=0,
\qquad \pi_{x'}\dot U(0)=0,
\qquad U(\theta)=0,
\]
and similarly for $V$.

\begin{proposition}
Let $U,V$ be piecewise smooth along $\gamma$ and satisfy these endpoint conditions. Then, in the distribution sense,
\begin{equation}
Q_\theta(U,V)=
\int_0^\theta
\left\langle
-\frac{D}{\partial t}\left(\frac{DU}{\partial t}\right)
+R(U,\dot\gamma)\dot\gamma,
V\right\rangle dt.
\tag{6.4$'$}
\end{equation}
\end{proposition}
The solutions of
\begin{equation}
-\frac{D}{\partial t}\left(\frac{DU}{\partial t}\right)
+R(U,\dot\gamma)\dot\gamma=0
\tag{6.5}
\end{equation}
are the Jacobi fields; they are precisely the infinitesimal variations of geodesics.

If $U$ is a Jacobi field on $[0,1]$ satisfying the initial conditions above, it is realized by a smooth variation through geodesics. For fixed $t_0$, the map from the admissible initial data to $U(t_0)$ is the differential of the geodesic-flow map
\[
\mathcal H_{t_0}:T_N^\perp\longrightarrow M,
\qquad (x,v)\longmapsto\exp_x(t_0v).
\]
We assume the minimality condition
\begin{equation}
\text{the length of }\gamma:[0,1]\to M
\text{ equals the distance from }\gamma(1)\text{ to }N.
\tag{6.6}
\end{equation}

\begin{proposition}
Assume (6.6). Then the differential of $\mathcal H_{t_0}$ at $(\gamma(0),\dot\gamma(0))$ is non-degenerate for every $t_0\in(0,1)$.
\end{proposition}
\begin{proof}
Suppose instead that it is degenerate for some $t_0$. Then there is a non-zero Jacobi field $U$ satisfying the admissible initial conditions and $U(t_0)=0$, hence $U$ lies in the kernel of $Q_{t_0}$. Since $\dot U(t_0)\ne0$, extend $U$ to a piecewise smooth field $\widetilde U$ on $[0,1]$ by setting it equal to zero on $[t_0,1]$. This is not a Jacobi field on $[0,1]$, so it is not in the kernel of $Q_1$. Choose $W$ with $Q_1(\widetilde U,W)\ne0$. Since $Q_1(\widetilde U,\widetilde U)=0$,
\[
Q_1(\widetilde U+\delta W,\widetilde U+\delta W)
=2\delta Q_1(\widetilde U,W)+\delta^2Q_1(W,W),
\]
which is negative for $\delta$ small with the appropriate sign. A variation having $\widetilde U+\delta W$ as variational field would then have vanishing first variation and negative second variation of the energy, contradicting minimality.
\end{proof}

Under the assumptions of the proposition one may choose an open neighborhood
$\mathcal U\subset T_N^\perp$ of
$\{(\gamma(0),\theta\dot\gamma(0));\ 0<\theta<1\}$ such that
\begin{equation}
(y,v)\in\mathcal U,
\quad0<\theta<1
\Longrightarrow (y,\theta v)\in\mathcal U,
\tag{6.7}
\end{equation}
and
\begin{equation}
\mathcal H_t|_{\mathcal U}
\quad\hbox{is a diffeomorphism onto }\mathcal H_t(\mathcal U).
\tag{6.8}
\end{equation}
This follows from the preceding proposition and the inverse function theorem.

\begin{theorem}
Under the assumptions of Proposition~6.3, $\mathcal U$ may be chosen so that
\[
\psi(x)=d(x,N)
\]
is $C^\infty$ on $\mathcal H_1(\mathcal U)$ and, for every
$x=\mathcal H_1(y,v)$ in that set, the distance $\psi(x)$ is realized, up to reparametrization, by a unique geodesic from $N$ to $x$, namely the geodesic $\gamma_{y,v}$.
\end{theorem}
\begin{proof}
Put $\widetilde\psi(x)=\norm v$ for $x=\mathcal H_1(y,v)$, $(y,v)\in\mathcal U$. This is a smooth function and is the length of the geodesic $\gamma_{y,v}$. If no sufficiently small $\mathcal U$ had the asserted property, there would be points $x_j=\mathcal H_1(y_j,v_j)$ converging to an interior point $\gamma(t_0)$ and distinct minimizing geodesics $\widetilde\gamma_j$ from $N$ to $x_j$. Compactness gives a limiting minimizing geodesic $\widetilde\gamma$ from $N$ to $\gamma(t_0)$ different from the corresponding segment of $\gamma$. The two terminal tangent vectors have the same norm but differ. Concatenating $\widetilde\gamma$ with $\gamma|_{[t_0,1]}$ therefore produces a minimizing curve with a corner, which is impossible. Hence $\widetilde\psi=d(\cdot,N)$ and the minimizing geodesic is unique.
\end{proof}

Return now to Section~5. Let $V\ge0$ and suppose $U_j=\{x_j\}$, $j=1,\ldots,N$, satisfy (5.14), without assuming (5.26)--(5.28). Apply the preceding discussion to
$M\setminus\{x_1,\ldots,x_N\}$ equipped with the Agmon metric $V(x)dx^2$. Let
\[
S_0=\inf_{j\ne k}d(x_j,x_k).
\]
For each $j$, let $\Omega_j$ consist of $x_j$ together with the interiors of all minimizing geodesics from $x_j$ to points of $M$ whose lengths are $<S_0$.

\begin{proposition}
The sets $\Omega_j$ are open and satisfy (5.11), (5.12), with $0$ replaced by $x_j$.
\end{proposition}
\begin{proof}
For $\delta>0$ small, a neighborhood of $B(x_j,\delta)$ has the desired properties by the local construction near a non-degenerate minimum. Put $N=\partial B(x_j,\delta)$. A minimizing geodesic from $x_j$ to a point $y\in B(x_j,S_0)\setminus B(x_j,\delta)$ decomposes into a minimizing segment from $x_j$ to $N$ and a minimizing segment from $N$ to $y$, and conversely. Since
\[
d(y,x_j)=\delta+d(y,N),
\]
Theorem~6.4 shows that $\Omega_j$ is open and that $d(\cdot,x_j)$ is smooth there. Reparametrizing the unique minimal geodesic as a null bicharacteristic of $q=\xi^2-V(x)$ gives (5.12).
\end{proof}

Choose the domains $M_j$ as in Section~2 and an interval $I(h)$ as in Section~5, so that all localized eigenvalues $\mu_\alpha$ are asymptotically equal; when more than one belongs to a given well, assume (5.29). If $d(U_j,U_k)=S_0$, choose a transverse hypersurface $\Gamma$ as in Remark~2.13. Whenever the local eigenfunctions admit asymptotic expansions near the intersection of $\Gamma$ with all minimizing geodesics joining the two wells, formulas (5.31)--(5.33) remain valid after decreasing $\Gamma$ if necessary. Hence the interaction coefficient has an asymptotic expansion in many cases even without (5.26), (5.28). For the lowest eigenvalues one obtains the following general bound.

\begin{theorem}
Choose $I(h)$ as in Theorem~5.9 and assume that $E_0$ can occur in the spectrum of the localized harmonic oscillators only as the lowest eigenvalue. If $\mu_\alpha,\mu_\beta\in I(h)$ are eigenvalues of $P_{M_{j(\alpha)}}$, $P_{M_{j(\beta)}}$ with
$d(U_{j(\alpha)},U_{j(\beta)})=S_0$, then there is a constant $C>0$ such that
\[
\frac1C h^{1/2}
<(-\widehat w_{\alpha,\beta})e^{S_0/h}
<C h^{1/2}.
\]
\end{theorem}

\begin{thebibliography}{99}
\bibitem{AbrahamMarsden}
R.~Abraham and J.~Marsden,
\emph{Foundations of Mechanics}, Benjamin/Cummings Publ. Company, 1978.

\bibitem{Agmon}
S.~Agmon,
\emph{Lectures on exponential decay of solutions of second-order elliptic equations},
Mathematical Notes 29, Princeton University Press.

\bibitem{CandelpergherNosmas}
B.~Candelpergher and J.-C.~Nosmas,
Propri\'et\'es spectrales d'op\'erateurs diff\'erentiels asymptotiques auto-adjoints,
Actes du Colloque de St Jean de Monts (Juin 1982), and Preprint.

\bibitem{CarmonaSimon}
R.~Carmona and B.~Simon,
Pointwise bounds on eigenfunctions and wave packets in $N$-body quantum systems V,
\emph{Comm. Math. Phys.} \textbf{80} (1981), 59--98.

\bibitem{Duistermaat}
J.~J.~Duistermaat,
Oscillatory integrals and unfolding of singularities,
\emph{C.P.A.M.} \textbf{27} (1974), 207--281.

\bibitem{GrigisSjostrand}
A.~Grigis and J.~Sj\"ostrand, In preparation.

\bibitem{HarrellDegeneracy}
E.~M.~Harrell,
On the rate of asymptotic eigenvalue degeneracy,
\emph{Comm. Math. Phys.} \textbf{60} (1978), 73--95.

\bibitem{HarrellDouble}
E.~M.~Harrell,
Double wells,
\emph{Comm. Math. Phys.} \textbf{75} (1980), 239--261.

\bibitem{HelfferRobert1}
B.~Helffer and D.~Robert,
Comportement semi-classique du spectre des hamiltoniens quantiques elliptiques,
\emph{Ann. Inst. Fourier} \textbf{31}, no.~3 (1981), 169--223.

\bibitem{HelfferRobert2}
B.~Helffer and D.~Robert,
Asymptotique compl\`ete des valeurs propres pour des \'equations diff\'erentielles sur $\R^n$,
\emph{Duke Math. J.} \textbf{49}, no.~4 (1982), 853--868.

\bibitem{HelfferRobert3}
B.~Helffer and D.~Robert,
Calcul fonctionnel par la transformation de Mellin et op\'erateurs admissibles,
\emph{J. Funct. Anal.} \textbf{53}, no.~3 (1983), 246--268.

\bibitem{HelfferRobert4}
B.~Helffer and D.~Robert,
Puits de potentiel et asymptotique semi-classique, Preprint.

\bibitem{JonaLasinio}
G.~Jona-Lasinio, F.~Martinelli and E.~Scoppola,
New approach to the semi-classical limit of quantum mechanics. I. Multiple tunneling in one dimension,
\emph{Comm. Math. Phys.} \textbf{80} (1981), 223--254.

\bibitem{LandauLifchitz}
L.~Landau and E.~Lifchitz,
\emph{M\'ecanique quantique, th\'eorie non relativiste}, Editions Mir, Moscou, 1966.

\bibitem{LascarSjostrand}
B.~Lascar and J.~Sj\"ostrand,
Equation de Schr\"odinger et propagation des singularit\'es pour des op\'erateurs pseudo-diff\'erentiels \`a caract\'eristiques r\'eelles de multiplicit\'e variable, II, Preprint.
(See also no.~I in \emph{Ast\'erisque} no.~95, 167--207.)

\bibitem{MansourMuller}
H.~M.~M.~Mansour and H.~J.~W.~M\"uller--Kristen,
Perturbation technik as an alternative to the W.K.B. method applied to the double well potential,
\emph{J. Math. Phys.} \textbf{23}(10), Oct. 1982.

\bibitem{Maslov}
V.~V.~Maslov,
\emph{Th\'eorie des perturbations et m\'ethodes asymptotiques}, Dunod.

\bibitem{Metivier}
G.~M\'etivier,
Analytic hypoellipticity for operators with multiple characteristics,
\emph{Comm. P.D.E.} \textbf{6}(1) (1981), 1--90.

\bibitem{Milnor}
J.~Milnor,
\emph{Morse Theory}, Princeton University Press, 1963.

\bibitem{ReedSimon}
M.~Reed and B.~Simon,
\emph{Methods of Modern Mathematical Physics}, Vol.~4, Academic Press, New York, 1978.

\bibitem{Persson}
A.~Persson,
Bounds for the discrete part of the spectrum of a semi-bounded Schr\"odinger operator,
\emph{Math. Scand.} \textbf{8} (1960), 143--153.

\bibitem{SimonLow1}
B.~Simon,
Semi-classical analysis of low lying eigenvalues I. Non-degenerate minima: Asymptotic expansions,
\emph{Ann. Inst. H. Poincar\'e} \textbf{38}, no.~3 (1983), 295--307 (and correction to appear).

\bibitem{SimonInstantons}
B.~Simon,
Instantons, double wells and large deviations,
\emph{Bull. Amer. Math. Soc.} \textbf{8}, no.~2 (March 1983), 323--326.
Complete proofs are given in: \emph{Semi-classical analysis of low lying eigenvalues II. Tunneling}, Preprint.

\bibitem{SjostrandSing}
J.~Sj\"ostrand,
\emph{Singularit\'es analytiques microlocales}, Ast\'erisque no.~95.

\bibitem{SjostrandWave}
J.~Sj\"ostrand,
Analytic wavefront sets and operators with multiple characteristics,
\emph{Hokkaido Math. J.} \textbf{12} (1983), 392--433.

\bibitem{Voros}
A.~Voros,
The return of the quartic oscillator,
\emph{Ann. Inst. H. Poincar\'e} \textbf{39}, no.~3 (1983), 211--338.

\bibitem{Davies}
E.~B.~Davies,
The twisting trick for double well Hamiltonians,
\emph{Comm. Math. Phys.} \textbf{85} (1982), 471--479.

\bibitem{DamburgPropin}
R.~J.~Damburg and R.~Kh.~Propin,
On one-dimensional model for exchange forces,
\emph{J. Chem. Phys.} \textbf{55}, no.~2 (1971), 612--616.

\bibitem{CombesDuclosSeiler}
J.~M.~Combes, P.~Duclos and R.~Seiler,
\begin{enumerate}[label=\arabic*.]
\item Krein's formula and one dimensional multiple wells,
\emph{J. Funct. Anal.} \textbf{52} (1983), 257--301.
\item Convergent expansions for tunneling,
\emph{Comm. Math. Phys.} \textbf{82} (1983), 229--245.
\end{enumerate}
\end{thebibliography}

\begin{flushright}
Received October 1983
\end{flushright}

\end{document}
